%% 3-4 グラフをかこう・読もう（第12回編集委員会：組 V3-2 が書いた）  % noruby
%% 流れ：お湯が冷める温度をはかる → ① 軸と目もり → ② 点となめらかな線（はてな？の答え）→ ③ x-t グラフと v-t グラフ  % noruby
%%       → ④ 内挿と外挿 → ⑤ グラフの形を読む（\R{右}{みぎ}\R{上}{あ}がり・いちばん\R{高}{たか}い所・交わる所）  % noruby
%% 第5巻 5-1・第3巻 3-5 の思い出そうが指す文「横軸を時間 t、たて軸を位置 x にしたグラフを、x-t グラフという」は変えないこと。  % noruby
\makeatletter
% とちゅうをはぶくしるし（1-18 の \bkqqnami と\R{同}{おな}じ形）  \bkthrenami  % noruby
\newcommand\bkthrenami{\tikz[baseline=-0.55ex, x=1mm, y=1mm]{\fill[white] (-0.3,-1.6) rectangle (5.1,1.6);\foreach \yy in {-0.7,0.7} {\draw[bkInk, line width=0.7pt, line cap=round] (0,\yy) .. controls (0.9,\yy+1) and (1.5,\yy-1) .. (2.4,\yy) .. controls (3.3,\yy+1) and (3.9,\yy-1) .. (4.8,\yy);}}}
% お湯の温度（ソラさん・ミナさんがはかった値）と、なめらかな線の式  % noruby
\def\bkthreSora{(0,80.0) (2,76.2) (4,71.9) (6,68.9) (8,65.1) (10,62.4) (12,59.2) (14,56.9) (16,54.1) (18,52.1) (20,49.7)}
\def\bkthreMina{(0,80.0) (2,75.7) (4,72.4) (6,68.4) (8,65.6) (10,62.0) (12,59.6) (14,56.5) (16,54.5) (18,51.7) (20,50.0)}
\def\bkthreT{20+60*exp(-0.035*x)}
% 説明の段の見出し（丸数字＋太字の1行）  \bkthrestep{番号}{見出し}  % noruby
\newcommand\bkthrestep[2]{\par\addvspace{6pt}\Needspace{4\baselineskip}\noindent
  {\bkcircled[bkOrange]{#1}\enspace\sffamily\bfseries\color{bkOrange!75!black}#2}\par\nobreak\vspace{2pt}}
% 例題の「しき」の行（折り返した行を字下げする。3-3 と\R{同}{おな}じ形）  % noruby
\newcommand\bkthreS{\shiki\hangindent=2.6\zw\hangafter=1\relax}
\newcommand\bkthreSi{\shiki\hangindent=5.1\zw\hangafter=1\relax}
\newcommand\bkthreI{\par\noindent\hangindent=5.1\zw\hangafter=1\relax\hspace*{2.6\zw}}
\newcommand\bkthreJ{\par\noindent\hangindent=5.1\zw\hangafter=0\relax}
\newcommand\bkthreC{\par\noindent\hangindent=2.6\zw\hangafter=0\relax}
\newcommand\bkthreK{\kotae\hangindent=2.6\zw\hangafter=1\relax}
\makeatother

\bkchapter{3}{4}{グラフをかこう・\R{読}{よ}もう}{\bklevel{1}}{グラフをかこう・\R{読}{よ}もう}
\bkmichishirube{1-17,1-18,2-3,2-6,3-2,3-3}{3-5,3-16,5-1,5-9}

\begin{goal}
\begin{itemize}[label={\color{\bkTT}●}, leftmargin=1.6\zw, itemsep=0pt, topsep=1pt]
\item \bklevel{1}\enspace \R{軸}{じく}に\R{量}{りょう}と\R{単位}{たん|い}を\R{書}{か}き、\R{目}{め}もりを\R{決}{き}めて、はかった\R{点}{てん}からグラフをかける。
\item \bklevel{1}\enspace $x$-$t$ グラフと $v$-$t$ グラフを\R{読}{よ}み\R{分}{わ}けられる。
\item \bklevel{1}\enspace グラフから\R{表}{ひょう}にない\R{値}{あたい}を\R{読}{よ}みとり、\R{線}{せん}を\R{先}{さき}へのばすときの\R{注意}{ちゅう|い}を\R{説明}{せつ|めい}できる。
\end{itemize}
\end{goal}

\begin{hatena}
カップのお\R{湯}{ゆ}が\R{冷}{さ}めていく\R{温度}{おん|ど}を2\R{分}{ふん}ごとにはかって、グラフに\R{点}{てん}をうちました。ソラさんは\R{点}{てん}を\R{順}{じゅん}にまっすぐな\R{線}{せん}で\R{結}{むす}び（おれ\R{線}{せん}）、ミナさんは\R{点}{てん}の\R{近}{ちか}くを\R{通}{とお}るなめらかな\R{線}{せん}を1\R{本}{ぽん}\R{引}{ひ}きました。どちらがよいのかな？　\R{表}{ひょう}にない5\R{分}{ふん}のときの\R{温度}{おん|ど}は、\R{何}{なん}℃くらいかな？
\begin{bkfigure}
\begin{tikzpicture}
\begin{axis}[bkplot, axis lines=left, name=A, width=54mm, height=36mm, xmin=0, xmax=21, ymin=42, ymax=83, xtick={0,5,10,15,20}, ytick={50,60,70,80}, xlabel={$t$（\R{分}{ふん}）}, ylabel={℃}]
  \addplot[bkOrange, line width=1.1pt] coordinates {\bkthreSora};
  \addplot[only marks, mark=*, mark size=1.4pt, bkInk] coordinates {\bkthreSora};
  \node[fill=white, inner sep=0.3pt] at (axis cs:0,45) {\bkthrenami};
\end{axis}
\node[font=\sffamily\bfseries\footnotesize, text=bkOrange!80!black, anchor=south] at ([yshift=5mm]A.north) {ソラさん：おれ\R{線}{せん}};
\begin{axis}[bkplot, axis lines=left, at={($(A.east)+(22mm,0)$)}, anchor=west, name=B, width=54mm, height=36mm, xmin=0, xmax=21, ymin=42, ymax=83, xtick={0,5,10,15,20}, ytick={50,60,70,80}, xlabel={$t$（\R{分}{ふん}）}, ylabel={℃}]
  \addplot[bkBlue, line width=1.3pt, domain=0:20.5] {\bkthreT};
  \addplot[only marks, mark=*, mark size=1.4pt, bkInk] coordinates {\bkthreSora};
  \node[fill=white, inner sep=0.3pt] at (axis cs:0,45) {\bkthrenami};
\end{axis}
\node[font=\sffamily\bfseries\footnotesize, text=bkBlue!80!black, anchor=south] at ([yshift=5mm]B.north) {ミナさん：なめらかな\R{線}{せん}};
\end{tikzpicture}
\end{bkfigure}
\end{hatena}

\clearpage
%% ---------------- やってみよう・思い出そう ----------------
\begin{yattemiyou}[お\R{湯}{ゆ}が\R{冷}{さ}めていく\R{温度}{おん|ど}をはかろう]
\noindent\begin{minipage}[c]{0.66\linewidth}
\begin{enumerate}[label={\bkcircled[bkGreen]{\arabic*}}, leftmargin=2.4\zw, itemsep=1pt, topsep=0pt]
\item おうちの\R{人}{ひと}といっしょに、カップに\R{熱}{あつ}いお\R{湯}{ゆ}を\R{入}{い}れる（やけどに\R{気}{き}をつけよう）。
\item \R{温度計}{おん|ど|けい}をお\R{湯}{ゆ}に\R{入}{い}れ、\R{時計}{と|けい}を\R{見}{み}て、0\R{分}{ふん}・2\R{分}{ふん}・4\R{分}{ぷん}…と、2\R{分}{ふん}ごとに20\R{分}{ぷん}まで\R{温度}{おん|ど}を\R{読}{よ}む（1-11）。
\item \R{読}{よ}んだ\R{温度}{おん|ど}を、\R{表}{ひょう}に\R{書}{か}いていく。
\end{enumerate}
\end{minipage}\hfill
\begin{minipage}[c]{0.3\linewidth}\centering
\begin{tikzpicture}[x=1mm, y=1mm]
  \bkcup{(12,0)}{20}{22}{0.75}{0}
  \fill[white, draw=bkGray, line width=0.7pt, rounded corners=1.2mm] (13,4) rectangle (17,34);
  \fill[bkRed] (15,6.5) circle (1.6);
  \fill[bkRed] (14.4,6.5) rectangle (15.6,24);
  \foreach \k in {0,...,8} \draw[bkInk, line width=0.35pt] (16,{10+\k*2.6}) -- (17,{10+\k*2.6});
  \bkclock*{(29,25)}{6}{3}{20}
\end{tikzpicture}
\end{minipage}
\par\vspace{4pt}
ソラさんがはかった\R{結果}{けっ|か}です（\R{部屋}{へ|や}の\R{温度}{おん|ど}は\qty{20}{\degreeCelsius}）。
\begin{bkcenter}
\begin{bktable}{colspec={Q[l,m,wd=15mm] *{11}{Q[c,m,wd=7.2mm]}}, rowsep=1.5pt, colsep=1pt, cells={font=\small}}
\R{時間}{じ|かん}（\R{分}{ふん}） & 0 & 2 & 4 & 6 & 8 & 10 & 12 & 14 & 16 & 18 & 20 \\
\R{温度}{おん|ど}（℃） & 80.0 & 76.2 & 71.9 & 68.9 & 65.1 & 62.4 & 59.2 & 56.9 & 54.1 & 52.1 & 49.7 \\
\end{bktable}
\end{bkcenter}
この\R{表}{ひょう}を、グラフにしてみよう。
\end{yattemiyou}

\noindent \R{温度}{おん|ど}は、0\R{分}{ふん}から20\R{分}{ぷん}までに、80.0℃から49.7℃まで\R{下}{さ}がりました。2\R{分}{ふん}ごとの\R{下}{さ}がり\R{方}{かた}を「あと−まえ」（2-3）で\R{計算}{けい|さん}すると、$76.2-80.0=-3.8$、$71.9-76.2=-4.3$、$68.9-71.9=-3.0$ …です。はじめは\R{速}{はや}く\R{下}{さ}がり、だんだんゆっくりになっていきます。ただ、$-3.8$、$-4.3$、$-3.0$ のように、\R{少}{すこ}しでこぼこもあります。このでこぼこを、グラフでどうあつかえばよいのでしょう。

\begin{omoidasou}
\bkrecall{1-17}{グラフをかく\R{前}{まえ}に、いちばん\R{大}{おお}きい\R{数}{かず}が\R{入}{はい}るように、1\R{目}{め}もりの\R{大}{おお}きさを\R{決}{き}める。}
\bkrecall{1-18}{おれ\R{線}{せん}グラフは、かわり\R{方}{かた}を\R{見}{み}るグラフ。かたむきが\R{急}{きゅう}なほど、はやくかわる。とちゅうをはぶくしるしは \bkthrenami。}
\bkrecall{2-3}{\R{速}{はや}さと\R{時間}{じ|かん}のグラフ（$v$-$t$ グラフ）。\R{速}{はや}さが\R{一定}{いっ|てい}なら、\R{長方形}{ちょう|ほう|けい}の\R{面積}{めん|せき}が\R{道}{みち}のり。}
\bkrecall{2-6}{はかった\R{数}{かず}には、ずれがある。\R{何回}{なん|かい}かはかって、ならして\R{考}{かんが}える。}
\bkrecall{3-2}{\R{点}{てん} $(x,\ y)$ は、\R{原点}{げん|てん}から\R{右}{みぎ}へ $x$、\R{上}{うえ}へ $y$ \R{進}{すす}んだ\R{所}{ところ}にかく。}
\bkrecall{3-3}{$y$ が $x$ の\R{関数}{かん|すう}のとき、\R{入}{い}れる\R{数}{かず} $x$ を\R{横軸}{よこ|じく}に、\R{出}{で}てくる\R{数}{かず} $y$ をたて\R{軸}{じく}にとる。}
\end{omoidasou}

%% ---------------- 数学の道具 ----------------
\bksection{doguu}{\R{数学}{すう|がく}の\R{道具}{どう|ぐ}}

\bksubsection{1}{\R{軸}{じく}と\R{目}{め}もりを\R{決}{き}める}

\leavevmode\label{bk:3-4:jiku}グラフは、\R{紙}{かみ}の\R{上}{うえ}にかいた「\R{表}{ひょう}の\R{絵}{え}」です。よい\R{絵}{え}にするには、かく\R{前}{まえ}に、\R{次}{つぎ}の①〜④を\R{決}{き}めます。

\begin{enumerate}[label={\bkcircled[bkTheme]{\arabic*}}, leftmargin=2.4\zw, itemsep=2pt, topsep=3pt]
\item \textbf{\R{横軸}{よこ|じく}とたて\R{軸}{じく}に、どの\R{量}{りょう}をとるか}\par\noindent \R{自分}{じ|ぶん}で\R{決}{き}めて\R{変}{か}えた\R{量}{りょう}や、\R{時間}{じ|かん}を\R{横軸}{よこ|じく}にとります。それにともなって\R{変}{か}わった\R{量}{りょう}（\R{結果}{けっ|か}）を、たて\R{軸}{じく}にとります。お\R{湯}{ゆ}の\R{実験}{じっ|けん}では、\R{時間}{じ|かん} $t$ を\R{横}{よこ}、\R{温度}{おん|ど} $T$ をたてにとります。\R{温度}{おん|ど}は\R{時間}{じ|かん}の\R{関数}{かん|すう}（3-3）で、$t$ が「\R{入}{い}れる\R{数}{かず}」、$T$ が「\R{出}{で}てくる\R{数}{かず}」だからです。
\item \textbf{\R{軸}{じく}に、\R{量}{りょう}の\R{名前}{な|まえ}と\R{単位}{たん|い}を\R{書}{か}く}\par\noindent 「$t$（\R{分}{ふん}）」「$T$（℃）」のように\R{書}{か}きます。\R{単位}{たん|い}がないと、10が10\R{秒}{びょう}なのか10\R{分}{ぷん}なのか、わかりません。
\item \textbf{1\R{目}{め}もりの\R{大}{おお}きさを\R{決}{き}める}\par\noindent いちばん\R{小}{ちい}さい\R{値}{あたい}から、いちばん\R{大}{おお}きい\R{値}{あたい}までが、\R{紙}{かみ}にできるだけ\R{大}{おお}きく\R{入}{はい}るようにします。1\R{目}{め}もりは、1・2・5、または10・20・50…や0.1・0.2・0.5…にします（1-17）。
\item \textbf{0から\R{始}{はじ}めるかを\R{決}{き}める}\par\noindent \R{値}{あたい}が0から\R{遠}{とお}い\R{所}{ところ}にかたまっているときは、とちゅうをはぶくしるし \bkthrenami を\R{軸}{じく}にかいて、0の\R{近}{ちか}くを\R{省}{はぶ}きます（1-18）。
\end{enumerate}

\textbf{なぜ1・2・5なのか}\quad \R{目}{め}もりは、あとで\R{読}{よ}むときに\R{楽}{らく}な\R{数}{かず}にします。1\R{目}{め}もりを3にすると、7は「2\R{目}{め}もりと、そのあと3\R{分}{ぶん}の1\R{目}{め}もり」の\R{所}{ところ}です。\R{目}{め}で3\R{等分}{とう|ぶん}するのは、むずかしいですね。1\R{目}{め}もりが2なら、7は「3\R{目}{め}もり\R{半}{はん}」、5なら「1\R{目}{め}もりと5\R{分}{ぶん}の2」で、ずっと\R{読}{よ}みやすくなります。1・2・5は、10をわりきれる\R{数}{かず}なので、10・20・50…と\R{大}{おお}きくしても、0.1・0.2・0.5…と\R{小}{ちい}さくしても、きりのよい\R{数}{かず}がならびます。

\begin{bkfigure}
\begin{tikzpicture}
\begin{axis}[bkplot, name=bad, width=52mm, height=48mm, xmin=0, xmax=63, ymin=0, ymax=330, xtick={0,21,42,63}, ytick={0,110,220,330}, xlabel={$t$（\R{分}{ふん}）}, ylabel={$T$（℃）}]
  \addplot[only marks, mark=*, mark size=1.2pt, bkInk] coordinates {\bkthreSora};
\end{axis}
\node[font=\sffamily\bfseries\footnotesize, text=bkRed!80!black, anchor=south] at ([yshift=6mm]bad.north) {×　\R{小}{ちい}さくて、\R{目}{め}もりも\R{読}{よ}みにくい};
\begin{axis}[bkplot, axis lines=left, at={($(bad.east)+(20mm,0)$)}, anchor=west, name=good, width=52mm, height=48mm, xmin=0, xmax=21, ymin=42, ymax=83, xtick={0,2,...,20}, xticklabels={0,,4,,8,,12,,16,,20}, ytick={50,55,...,80}, xlabel={$t$（\R{分}{ふん}）}, ylabel={$T$（℃）}]
  \addplot[only marks, mark=*, mark size=1.5pt, bkInk] coordinates {\bkthreSora};
  \node[fill=white, inner sep=0.3pt] at (axis cs:0,45) {\bkthrenami};
\end{axis}
\node[font=\sffamily\bfseries\footnotesize, text=bkGreen!50!black, anchor=south] at ([yshift=6mm]good.north) {○　\R{大}{おお}きくて、\R{読}{よ}みやすい};
\end{tikzpicture}
\bkfigcaption{\R{同}{おな}じ\R{表}{ひょう}から\R{作}{つく}った2つのグラフ。\R{左}{ひだり}は1\R{目}{め}もりが21\R{分}{ぷん}・110℃、\R{右}{みぎ}は2\R{分}{ふん}・5℃。}
\end{bkfigure}

\R{左}{ひだり}のグラフは、0℃から330℃までの\R{目}{め}もりなので、\R{点}{てん}が\R{下}{した}のほうに\R{小}{ちい}さくかたまり、\R{変}{か}わり\R{方}{かた}が\R{見}{み}えません。1\R{目}{め}もりも21や110で、\R{読}{よ}みにくい\R{数}{かず}です。\R{右}{みぎ}のグラフは、\R{温度}{おん|ど}の\R{目}{め}もりを45℃より\R{下}{した}で\R{省}{はぶ}き、1\R{目}{め}もりを5℃にしたので、\R{点}{てん}の\R{並}{なら}び\R{方}{かた}がよく\R{見}{み}えます。

\bksubsection{2}{\R{点}{てん}をうって、なめらかな\R{線}{せん}を\R{引}{ひ}く}

\leavevmode\label{bk:3-4:sen}\R{点}{てん}をうったら、\R{線}{せん}を\R{引}{ひ}きます。\bkhatenaref のように、\R{引}{ひ}き\R{方}{かた}は2つ\R{考}{かんが}えられます。おれ\R{線}{せん}にするか、なめらかな\R{線}{せん}にするかです。どちらがよいかを\R{考}{かんが}えるために、\R{同}{おな}じカップのお\R{湯}{ゆ}を、ソラさんとミナさんが、\R{同時}{どう|じ}にはかったとしましょう。

\begin{bkcenter}
\begin{bktable}{colspec={Q[l,m,wd=15mm] *{11}{Q[c,m,wd=7.2mm]}}, rowsep=1.5pt, colsep=1pt, cells={font=\small}}
\R{時間}{じ|かん}（\R{分}{ふん}） & 0 & 2 & 4 & 6 & 8 & 10 & 12 & 14 & 16 & 18 & 20 \\
ソラ（℃） & 80.0 & 76.2 & 71.9 & 68.9 & 65.1 & 62.4 & 59.2 & 56.9 & 54.1 & 52.1 & 49.7 \\
ミナ（℃） & 80.0 & 75.7 & 72.4 & 68.4 & 65.6 & 62.0 & 59.6 & 56.5 & 54.5 & 51.7 & 50.0 \\
\end{bktable}
\bktabcaption{\R{同}{おな}じお\R{湯}{ゆ}を、\R{2人}{ふたり}で\R{同時}{どう|じ}にはかった。\R{2人}{ふたり}の\R{値}{あたい}は、\R{少}{すこ}しずつちがう。}
\end{bkcenter}

\begin{bkfigure}
\begin{tikzpicture}
\begin{axis}[bkplot, axis lines=left, name=A, width=54mm, height=48mm, xmin=0, xmax=21, ymin=42, ymax=83, xtick={0,5,10,15,20}, ytick={50,60,70,80}, xlabel={$t$（\R{分}{ふん}）}, ylabel={$T$（℃）}]
  \addplot[bkOrange, line width=1pt] coordinates {\bkthreSora};
  \addplot[bkGreen!70!black, line width=1pt] coordinates {\bkthreMina};
  \addplot[only marks, mark=*, mark size=1.3pt, bkOrange] coordinates {\bkthreSora};
  \addplot[only marks, mark=square*, mark size=1.3pt, bkGreen!70!black] coordinates {\bkthreMina};
  \node[fill=white, inner sep=0.3pt] at (axis cs:0,45) {\bkthrenami};
\end{axis}
\node[font=\sffamily\bfseries\footnotesize, anchor=south] at ([yshift=6mm]A.north) {おれ\R{線}{せん}：\R{2人}{ふたり}の\R{線}{せん}がちがう};
\begin{axis}[bkplot, axis lines=left, at={($(A.east)+(20mm,0)$)}, anchor=west, name=B, width=54mm, height=48mm, xmin=0, xmax=21, ymin=42, ymax=83, xtick={0,5,10,15,20}, ytick={50,60,70,80}, xlabel={$t$（\R{分}{ふん}）}, ylabel={$T$（℃）}]
  \addplot[bkBlue, line width=1.3pt, domain=0:20.5] {\bkthreT};
  \addplot[only marks, mark=*, mark size=1.3pt, bkOrange] coordinates {\bkthreSora};
  \addplot[only marks, mark=square*, mark size=1.3pt, bkGreen!70!black] coordinates {\bkthreMina};
  \node[fill=white, inner sep=0.3pt] at (axis cs:0,45) {\bkthrenami};
\end{axis}
\node[font=\sffamily\bfseries\footnotesize, anchor=south] at ([yshift=6mm]B.north) {なめらかな\R{線}{せん}：1\R{本}{ぽん}で\R{2人}{ふたり}に\R{合}{あ}う};
\end{tikzpicture}
\bkfigcaption{●がソラさん、■がミナさん。}
\end{bkfigure}

\textbf{\R{見方}{み|かた}1：\R{2人}{ふたり}の\R{線}{せん}をくらべる}\quad おれ\R{線}{せん}でつなぐと、ソラさんの\R{線}{せん}とミナさんの\R{線}{せん}は、ぎざぎざの\R{形}{かたち}がちがい、\R{何回}{なん|かい}も\R{交}{まじ}わります。\R{同}{おな}じお\R{湯}{ゆ}をはかったのに、\R{線}{せん}の\R{形}{かたち}がちがうのは、おかしいですね。ぎざぎざは、お\R{湯}{ゆ}の\R{冷}{さ}め\R{方}{かた}ではなく、\R{温度計}{おん|ど|けい}を\R{読}{よ}むときのずれ（2-6）から\R{生}{う}まれたものだからです。なめらかな\R{線}{せん}なら、1\R{本}{ぽん}で\R{2人}{ふたり}の\R{点}{てん}にどちらも\R{合}{あ}います。

\textbf{\R{見方}{み|かた}2：\R{数}{かず}で\R{考}{かんが}える}\quad 6\R{分}{ぷん}のとき、ソラさんは68.9℃、ミナさんは68.4℃でした。どちらかが\R{正}{ただ}しいのではなく、\R{本当}{ほん|とう}の\R{温度}{おん|ど}はこの\R{近}{ちか}くのどこかです。4\R{分}{ぷん}と8\R{分}{ぷん}の\R{値}{あたい}から\R{考}{かんが}えると、6\R{分}{ぷん}の\R{本当}{ほん|とう}の\R{温度}{おん|ど}は、\R{2人}{ふたり}の\R{値}{あたい}のあいだの\R{約}{やく}\qty{68.6}{\degreeCelsius}だと\R{考}{かんが}えられます。1つ1つの\R{点}{てん}は\R{少}{すこ}しずれているので、\R{全部}{ぜん|ぶ}の\R{点}{てん}の「まん\R{中}{なか}」を\R{通}{とお}る\R{線}{せん}のほうが、\R{本当}{ほん|とう}の\R{変}{か}わり\R{方}{かた}に\R{近}{ちか}いのです。

\begin{taisetsu}
\R{実験}{じっ|けん}で、あいだも\R{続}{つづ}けて\R{変}{か}わる\R{量}{りょう}（\R{時間}{じ|かん}・\R{温度}{おん|ど}・\R{長}{なが}さなど）をはかったときは、\R{点}{てん}と\R{点}{てん}を\R{結}{むす}ぶおれ\R{線}{せん}ではなく、\R{点}{てん}の\R{近}{ちか}くを\R{通}{とお}るなめらかな\R{線}{せん}を1\R{本}{ぽん}\R{引}{ひ}く。\R{点}{てん}がまっすぐに\R{並}{なら}ぶなら\R{定規}{じょう|ぎ}で\R{直線}{ちょく|せん}を、\R{曲}{ま}がって\R{並}{なら}ぶならフリーハンドで\R{曲線}{きょく|せん}を\R{引}{ひ}く。\R{点}{てん}が\R{線}{せん}の\R{上}{うえ}と\R{下}{した}に\R{同}{おな}じくらい\R{散}{ち}らばるようにする。
\end{taisetsu}

\bkhatenaref の\R{答}{こた}え：ミナさんのなめらかな\R{線}{せん}がよいのです。5\R{分}{ふん}のときの\R{温度}{おん|ど}は、\R{線}{せん}から\R{読}{よ}むと\R{約}{やく}\qty{70}{\degreeCelsius}です（\R{読}{よ}み\R{方}{かた}は④でくわしく）。

\textbf{おれ\R{線}{せん}を\R{使}{つか}うとき}\quad 1-18 では、1\R{時間}{じ|かん}ごとの\R{気温}{き|おん}を、おれ\R{線}{せん}グラフにしました。おれ\R{線}{せん}は、「いくつかの\R{時刻}{じ|こく}の\R{記録}{き|ろく}」をならべて、\R{変}{か}わり\R{方}{かた}をひと\R{目}{め}で\R{見}{み}るためのグラフです。\R{月}{つき}ごとの\R{電気代}{でん|き|だい}や、\R{日}{にち}ごとの\R{最高気温}{さい|こう|き|おん}のように、あいだの\R{値}{あたい}がないものや、あいだの\R{変}{か}わり\R{方}{かた}がわからないものは、おれ\R{線}{せん}で\R{表}{あらわ}します。

\textbf{まっすぐに\R{並}{なら}ぶとき}\quad \R{電熱線}{でん|ねつ|せん}にかける\R{電圧}{でん|あつ}を\R{変}{か}えて、\R{電流}{でん|りゅう}をはかりました（2-15）。\R{点}{てん}は、ほぼまっすぐに\R{並}{なら}び、0Vのときは\R{電流}{でん|りゅう}も0です。だから、\R{定規}{じょう|ぎ}を\R{当}{あ}てて、\R{原点}{げん|てん}を\R{通}{とお}り、\R{点}{てん}の\R{近}{ちか}くを\R{通}{とお}る\R{直線}{ちょく|せん}を\R{引}{ひ}きます。
\begin{bkfigure}
\begin{tikzpicture}
\begin{axis}[bkplot, width=70mm, height=46mm, xmin=0, xmax=6.8, ymin=0, ymax=0.142, xtick={0,1.5,3,4.5,6}, ytick={0,0.03,0.06,0.09,0.12},
  yticklabel style={/pgf/number format/fixed, /pgf/number format/precision=2}, xticklabel style={/pgf/number format/fixed}, xlabel={$V$（V）}, ylabel={$I$（A）}]
  \addplot[bkBlue, line width=1.3pt, domain=0:6.5] {x/50};
  \addplot[only marks, mark=*, mark size=1.5pt, bkInk] coordinates {(0,0) (1.5,0.031) (3,0.059) (4.5,0.091) (6,0.119)};
\end{axis}
\end{tikzpicture}
\bkfigcaption{はかった\R{点}{てん}（\R{黒}{くろ}）と、\R{定規}{じょう|ぎ}で\R{引}{ひ}いた\R{直線}{ちょく|せん}（\R{青}{あお}）。}
\end{bkfigure}

\bksubsection{3}{\R{位置}{い|ち}と\R{時間}{じ|かん}のグラフ：$x$-$t$ グラフ}

\leavevmode\label{bk:3-4:xt}ソラさんが、\R{家}{いえ}から\qty{700}{m}\R{先}{さき}の\R{図書館}{と|しょ|かん}まで\R{行}{い}きました。\R{家}{いえ}からの\R{道}{みち}を\R{数直線}{すう|ちょく|せん}にして（2-2）、\R{家}{いえ}を0、\R{図書館}{と|しょ|かん}の\R{向}{む}きを＋とし、1\R{分}{ぷん}ごとにソラさんの\R{位置}{い|ち} $x$〔m〕を\R{書}{か}きます。はじめの5\R{分}{ふん}は1\R{分}{ぷん}に\qty{60}{m}ずつ\R{歩}{ある}き、\R{次}{つぎ}の3\R{分}{ぷん}は\R{信号}{しん|ごう}で\R{止}{と}まり、\R{最後}{さい|ご}の4\R{分}{ぷん}は1\R{分}{ぷん}に\qty{100}{m}ずつ\R{急}{いそ}いで\R{歩}{ある}きました。

\begin{bkcenter}
\begin{bktable}{colspec={Q[l,m,wd=16mm] *{13}{Q[c,m,wd=5.6mm]}}, rowsep=1.5pt, colsep=1.4pt, cells={font=\small}}
$t$（\R{分}{ふん}） & 0 & 1 & 2 & 3 & 4 & 5 & 6 & 7 & 8 & 9 & 10 & 11 & 12 \\
$x$（m） & 0 & 60 & 120 & 180 & 240 & 300 & 300 & 300 & 300 & 400 & 500 & 600 & 700 \\
\end{bktable}
\bktabcaption{ソラさんの\R{位置}{い|ち}。5\R{分}{ふん}から8\R{分}{ぷん}までは、300\,mの\R{所}{ところ}で\R{止}{と}まっている。}
\end{bkcenter}

\begin{bkfigure}
\begin{tikzpicture}[x=1mm, y=1mm]
  % 左：道をたてにした数直線（1 m＝0.08 mm）。右：x-t グラフ（1分＝7 mm）
  \draw[-{Stealth[length=4pt]}, bkInk, line width=0.8pt] (6,0) -- (6,62);
  \foreach \m in {0,100,...,700} {\draw[bkInk, line width=0.5pt] (4.8,\m*0.08) -- (7.2,\m*0.08);}
  \foreach \m in {0,200,400,600} {\node[font=\footnotesize, anchor=east] at (4,\m*0.08) {\m};}
  \node[font=\sffamily\footnotesize, anchor=south] at (6,63) {\R{道}{みち}（m）};
  \node[font=\sffamily\footnotesize, anchor=west, fill=white, inner sep=0.5pt] at (8,0) {\R{家}{いえ}};
  \node[font=\sffamily\footnotesize, anchor=west] at (8,56) {\R{図書館}{と|しょ|かん}};
  \begin{scope}[shift={(32,0)}]
  \foreach \k in {0,...,12} {\draw[bkLine, line width=0.3pt] (\k*6.5,0) -- (\k*6.5,58);}
  \foreach \m in {100,200,...,700} {\draw[bkLine, line width=0.3pt] (0,\m*0.08) -- (80,\m*0.08);}
  \draw[-{Stealth[length=4pt]}, bkInk, line width=0.8pt] (0,0) -- (83,0);
  \draw[-{Stealth[length=4pt]}, bkInk, line width=0.8pt] (0,0) -- (0,62);
  \foreach \k in {0,2,...,12} {\node[font=\footnotesize, anchor=north] at (\k*6.5,-0.8) {\k};}
  \foreach \m in {0,200,400,600} {\node[font=\footnotesize, anchor=east] at (-0.8,\m*0.08) {\m};}
  \node[font=\sffamily\footnotesize, anchor=west] at (84,0) {$t$（\R{分}{ふん}）};
  \node[font=\sffamily\footnotesize, anchor=south] at (0,63) {$x$（m）};
  \draw[bkBlue, line width=1.4pt] (0,0) -- (32.5,24) -- (52,24) -- (78,56);
  \foreach \t/\x in {0/0,1/60,2/120,3/180,4/240,5/300,6/300,7/300,8/300,9/400,10/500,11/600,12/700} {\fill[bkInk] (\t*6.5,\x*0.08) circle (0.8);}
  \end{scope}
  \foreach \x in {180,300,600} {\draw[bkGray, line width=0.5pt, dash pattern=on 1.5pt off 1.2pt] (7.5,\x*0.08) -- (21,\x*0.08);}
  \fill[bkOrange] (6,14.4) circle (1.1) (6,24) circle (1.1) (6,48) circle (1.1);
\end{tikzpicture}
\bkfigcaption{\R{道}{みち}の\R{数直線}{すう|ちょく|せん}をたてに\R{立}{た}てて、\R{時刻}{じ|こく}ごとに\R{右}{みぎ}へずらして\R{点}{てん}をうつと、$x$-$t$ グラフになる。}
\end{bkfigure}

\R{道}{みち}の\R{数直線}{すう|ちょく|せん}だけでは、ソラさんがどこにいたかはわかりますが、いつそこにいたかはわかりません。そこで、\R{数直線}{すう|ちょく|せん}をたてに\R{立}{た}てて、1\R{分}{ぷん}ごとに\R{右}{みぎ}へずらしながら、その\R{時刻}{じ|こく}の\R{位置}{い|ち}に\R{点}{てん}をうちます。3\R{分}{ぷん}のときは180\,m、5\R{分}{ふん}のときは300\,mの\R{高}{たか}さです。\R{点}{てん}を\R{線}{せん}で\R{結}{むす}ぶと、$x$-$t$ グラフになります。\R{歩}{ある}いているあいだ、\R{位置}{い|ち}は\R{時間}{じ|かん}とともに\R{続}{つづ}けて\R{変}{か}わるので、\R{点}{てん}と\R{点}{てん}のあいだも\R{線}{せん}でつなげます。

このように、\textbf{\R{横軸}{よこ|じく}を\R{時間}{じ|かん} $t$、たて\R{軸}{じく}を\R{位置}{い|ち} $x$ にしたグラフを、$x$-$t$ グラフという}。\bkidx{えっくすてぃーぐらふ}{$x$-$t$ グラフ}\R{位置}{い|ち} $x$ は\R{時刻}{じ|こく} $t$ の\R{関数}{かん|すう}なので、3-3 の\R{書}{か}き\R{方}{かた}で $x(3)=180$、$x(5)=300$ と\R{書}{か}けます。

\textbf{$x$-$t$ グラフの\R{読}{よ}み\R{方}{かた}}
\begin{itemize}[leftmargin=1.6\zw, itemsep=1pt, topsep=2pt, label={\color{\bkTT}●}]
\item \textbf{\R{線}{せん}の\R{高}{たか}さ}は、その\R{時刻}{じ|こく}にいる\R{場所}{ば|しょ}。12\R{分}{ふん}の\R{高}{たか}さ700は、\R{図書館}{と|しょ|かん}にいること。
\item \textbf{\R{線}{せん}が\R{水平}{すい|へい}}な\R{所}{ところ}は、\R{時間}{じ|かん}がたっても\R{位置}{い|ち}が\R{変}{か}わらない、つまり\textbf{\R{止}{と}まっている}。5\R{分}{ふん}から8\R{分}{ぷん}。
\item \textbf{\R{線}{せん}が\R{急}{きゅう}}な\R{所}{ところ}ほど、1\R{分}{ぷん}に\R{進}{すす}む\R{道}{みち}のりが\R{大}{おお}きく、\textbf{\R{速}{はや}い}（1-18 の「かたむきが\R{急}{きゅう}なほど、はやくかわる」）。8\R{分}{ぷん}から12\R{分}{ふん}のほうが、0\R{分}{ふん}から5\R{分}{ふん}より\R{急}{きゅう}。
\item もし\R{家}{いえ}へもどれば、\R{位置}{い|ち}が\R{小}{ちい}さくなるので、\R{線}{せん}は\textbf{\R{右}{みぎ}\R{下}{さ}がり}になる。
\end{itemize}

\textbf{$v$-$t$ グラフとならべる}\quad \R{同}{おな}じ\R{道}{みち}のりを、\R{速}{はや}さ $v$ と\R{時間}{じ|かん} $t$ のグラフ（$v$-$t$ グラフ、2-3）にもかいてみます。0\R{分}{ふん}から5\R{分}{ふん}は\R{速}{はや}さ\qty{60}{m/min}、5\R{分}{ふん}から8\R{分}{ぷん}は0、8\R{分}{ぷん}から12\R{分}{ふん}は\qty{100}{m/min}です。
\begin{bkfigure}
\begin{tikzpicture}
\begin{axis}[bkplot, name=xt, width=104mm, height=44mm, xmin=0, xmax=13, ymin=0, ymax=800, xtick={0,1,...,12}, ytick={0,200,400,600}, xlabel={$t$（\R{分}{ふん}）}, ylabel={$x$（m）}]
  \addplot[bkBlue, line width=1.4pt] coordinates {(0,0) (5,300) (8,300) (12,700)};
  \node[font=\sffamily\footnotesize, anchor=south, text=bkBlue!80!black] at (axis cs:6.5,380) {\R{水平}{すい|へい}＝\R{止}{と}まっている};
\end{axis}
\begin{axis}[bkplot, at={($(xt.south)+(0,-12mm)$)}, anchor=north, width=104mm, height=36mm, xmin=0, xmax=13, ymin=0, ymax=130, xtick={0,1,...,12}, ytick={0,60,100}, xlabel={$t$（\R{分}{ふん}）}, ylabel={$v$（m/min）}]
  \addplot[fill=bkBlue!15, draw=none] coordinates {(0,0) (0,60) (5,60) (5,0)} \closedcycle;
  \addplot[fill=bkBlue!15, draw=none] coordinates {(8,0) (8,100) (12,100) (12,0)} \closedcycle;
  \addplot[bkOrange, line width=1.4pt] coordinates {(0,60) (5,60)};
  \addplot[bkOrange, line width=1.4pt] coordinates {(5,0) (8,0)};
  \addplot[bkOrange, line width=1.4pt] coordinates {(8,100) (12,100)};
  \node[font=\sffamily\footnotesize] at (axis cs:2.5,30) {$60\times5=300$};
  \node[font=\sffamily\footnotesize] at (axis cs:10,50) {$100\times4=400$};
\end{axis}
\end{tikzpicture}
\bkfigcaption{\R{上}{うえ}が $x$-$t$ グラフ、\R{下}{した}が $v$-$t$ グラフ。\R{同}{おな}じ\R{道}{みち}のりを2つのグラフで\R{表}{あらわ}した。}
\end{bkfigure}

\begin{bkcenter}
\begin{bktable}{colspec={Q[l,m,wd=30mm] Q[l,m,wd=44mm] Q[l,m,wd=44mm]}, rowsep=1.5pt, cells={font=\small}, row{1}={font=\sffamily\bfseries\small}}
 & $x$-$t$ グラフ & $v$-$t$ グラフ \\
たて\R{軸}{じく} & \R{位置}{い|ち}（どこにいるか） & \R{速}{はや}さ（どれだけ\R{速}{はや}いか） \\
\R{止}{と}まっているとき & \R{水平}{すい|へい}な\R{線}{せん}（\R{高}{たか}さはその\R{場所}{ば|しょ}） & \R{高}{たか}さ0の\R{線}{せん} \\
\R{速}{はや}く\R{進}{すす}むとき & \R{急}{きゅう}な\R{線}{せん} & \R{高}{たか}い\R{線}{せん} \\
\R{進}{すす}んだ\R{道}{みち}のり & たての\R{変化量}{へん|か|りょう}（あと−まえ） & \R{線}{せん}の\R{下}{した}の\R{面積}{めん|せき} \\
\end{bktable}
\bktabcaption{$x$-$t$ グラフと $v$-$t$ グラフのちがい。}
\end{bkcenter}
\R{進}{すす}んだ\R{道}{みち}のりを2つのグラフでくらべます。$x$-$t$ グラフでは、たての\R{変化量}{へん|か|りょう}で $700-0=700$（m）。$v$-$t$ グラフでは、\R{長方形}{ちょう|ほう|けい}の\R{面積}{めん|せき}（2-3）で $60\times5+100\times4=300+400=700$（m）。\R{同}{おな}じになります。

\bksubsection{4}{\R{表}{ひょう}にない\R{値}{あたい}を\R{読}{よ}む：\R{内挿}{ない|そう}と\R{外挿}{がい|そう}}

\leavevmode\label{bk:3-4:naisou}なめらかな\R{線}{せん}を\R{引}{ひ}くと、はかっていない\R{時刻}{じ|こく}の\R{温度}{おん|ど}も\R{読}{よ}みとれます。

\begin{bkfigure}
\begin{tikzpicture}
\begin{axis}[bkplot, axis lines=left, width=100mm, height=62mm, xmin=0, xmax=21, ymin=42, ymax=83, xtick={0,2,...,20}, ytick={50,55,...,80}, minor xtick={1,3,...,19}, minor ytick={46,47,...,82}, minor grid style={bkLine!60, line width=0.15pt}, grid=both, xlabel={$t$（\R{分}{ふん}）}, ylabel={$T$（℃）}]
  \addplot[bkBlue, line width=1.3pt, domain=0:20.5] {\bkthreT};
  \addplot[only marks, mark=*, mark size=1.5pt, bkInk] coordinates {\bkthreSora};
  \draw[bkRed, line width=0.8pt, dash pattern=on 2.5pt off 1.5pt, -{Stealth[length=4pt]}] (axis cs:5,45) -- (axis cs:5,70.4) -- (axis cs:0.3,70.4);
  \draw[bkGreen!60!black, line width=0.8pt, dash pattern=on 2.5pt off 1.5pt, -{Stealth[length=4pt]}] (axis cs:0,65) -- (axis cs:8.2,65) -- (axis cs:8.2,45.4);
  \node[font=\sffamily\footnotesize, text=bkRed!80!black, anchor=south west, fill=white, inner sep=0.8pt] at (axis cs:5.3,71.2) {① 5\R{分}{ふん}→\R{約}{やく}70℃};
  \node[font=\sffamily\footnotesize, text=bkGreen!40!black, anchor=east, fill=white, inner sep=0.8pt] at (axis cs:7.4,57) {② 65℃→\R{約}{やく}8\R{分}{ぷん}};
  \node[fill=white, inner sep=0.3pt] at (axis cs:0,45) {\bkthrenami};
\end{axis}
\end{tikzpicture}
\bkfigcaption{\R{線}{せん}から\R{読}{よ}む。①は\R{横軸}{よこ|じく}から\R{上}{うえ}へ、②はたて\R{軸}{じく}から\R{横}{よこ}へ\R{進}{すす}む。}
\end{bkfigure}

\bkthrestep{1}{\R{横軸}{よこ|じく}の\R{値}{あたい}から、たての\R{値}{あたい}を\R{読}{よ}む}
5\R{分}{ふん}のときの\R{温度}{おん|ど}を\R{読}{よ}みます。\R{横軸}{よこ|じく}の5から\R{真上}{ま|うえ}へ\R{進}{すす}み、\R{線}{せん}に\R{当}{あ}たった\R{所}{ところ}から\R{左}{ひだり}へ\R{進}{すす}んで、たて\R{軸}{じく}の\R{目}{め}もりを\R{読}{よ}みます。\R{約}{やく}\qty{70}{\degreeCelsius}です。

\textbf{\R{数}{かず}でも\R{確}{たし}かめる}\quad 5\R{分}{ふん}は、4\R{分}{ぷん}と6\R{分}{ぷん}のちょうどまん\R{中}{なか}です。4\R{分}{ぷん}は71.9℃、6\R{分}{ぷん}は68.9℃なので、そのまん\R{中}{なか}は $(71.9+68.9)\div2=140.8\div2=70.4$ で、\R{約}{やく}70.4℃です。グラフで\R{読}{よ}んだ\R{値}{あたい}と\R{合}{あ}います。

\bkthrestep{2}{たての\R{値}{あたい}から、\R{横軸}{よこ|じく}の\R{値}{あたい}を\R{読}{よ}む}
\R{反対}{はん|たい}に、\R{温度}{おん|ど}が\qty{65}{\degreeCelsius}になったのは、\R{何分}{なん|ぷん}のときでしょう。たて\R{軸}{じく}の65から\R{横}{よこ}へ\R{進}{すす}み、\R{線}{せん}に\R{当}{あ}たった\R{所}{ところ}から\R{下}{した}へ\R{進}{すす}むと、\R{約}{やく}8\R{分}{ぷん}です。8\R{分}{ぷん}ではかった65.1℃とも\R{合}{あ}います。

このように、はかった\R{点}{てん}と\R{点}{てん}のあいだの\R{値}{あたい}を、グラフの\R{線}{せん}から\R{読}{よ}みとることを\textbf{\R{内挿}{ない|そう}}といいます。\bkidx{ないそう}{\R{内挿}{ない|そう}}\R{内側}{うち|がわ}（はかった\R{範囲}{はん|い}の\R{中}{なか}）を「おぎなう」という\R{意味}{い|み}です。

\bkthrestep{3}{\R{線}{せん}を\R{先}{さき}へのばすのは\R{危}{あぶ}ない：\R{外挿}{がい|そう}}
では、40\R{分}{ぷん}のときの\R{温度}{おん|ど}を\R{予想}{よ|そう}してみましょう。はじめの2\R{分}{ふん}で\qty{3.8}{\degreeCelsius}、1\R{分}{ぷん}あたり $3.8\div2=1.9$（℃）\R{下}{さ}がりました。このまま\R{同}{おな}じ\R{速}{はや}さで\R{下}{さ}がり\R{続}{つづ}けると\R{考}{かんが}えて、まっすぐな\R{線}{せん}をのばすと、40\R{分}{ぷん}では
\[ 80.0-1.9\times40=80.0-76=4.0 \]
で、\qty{4}{\degreeCelsius}になってしまいます。でも、\R{部屋}{へ|や}の\R{温度}{おん|ど}は\qty{20}{\degreeCelsius}です。お\R{湯}{ゆ}が、\R{部屋}{へ|や}より\R{冷}{つめ}たくなることはありません。

\begin{bkfigure}
\begin{tikzpicture}
\begin{axis}[bkplot, width=104mm, height=60mm, xmin=0, xmax=64, ymin=0, ymax=88, xtick={0,10,...,60}, ytick={0,20,40,60,80}, xlabel={$t$（\R{分}{ふん}）}, ylabel={$T$（℃）}]
  \addplot[bkGray, line width=0.8pt, dash pattern=on 3pt off 2pt, domain=0:62] {20};
  \node[font=\sffamily\footnotesize, text=bkGray!80!black, anchor=north east] at (axis cs:62,18.5) {\R{部屋}{へ|や}の\R{温度}{おん|ど} 20℃};
  \addplot[bkRed, line width=1.1pt, dash pattern=on 4pt off 2pt, domain=0:42] {80-1.9*x};
  \addplot[bkBlue, line width=1.3pt, domain=0:20] {\bkthreT};
  \addplot[bkBlue!60, line width=1.1pt, dash pattern=on 4pt off 2pt, domain=20:62] {\bkthreT};
  \addplot[only marks, mark=*, mark size=1.4pt, bkInk] coordinates {\bkthreSora};
  \addplot[only marks, mark=o, mark size=1.8pt, bkInk, line width=0.8pt] coordinates {(30,41.0) (40,34.8) (50,30.4) (60,27.3)};
  \node[font=\sffamily\footnotesize, text=bkRed!80!black, anchor=south east, align=right, inner sep=0.8pt] at (axis cs:28,1.5) {まっすぐのばすと\\40\R{分}{ぷん}で4℃？};
  \node[font=\sffamily\footnotesize, anchor=south west, fill=white, inner sep=0.8pt] at (axis cs:34,40) {○はあとではかった\R{値}{あたい}};
\end{axis}
\end{tikzpicture}
\bkfigcaption{\R{赤}{あか}の\R{点線}{てん|せん}は、はじめの\R{下}{さ}がり\R{方}{かた}のまま、まっすぐにのばした\R{線}{せん}。\R{本当}{ほん|とう}の\R{温度}{おん|ど}（○）は、\R{部屋}{へ|や}の\R{温度}{おん|ど}に\R{近}{ちか}づいていく。}
\end{bkfigure}

あとで\R{本当}{ほん|とう}にはかると、30\R{分}{ぷん}で41.0℃、40\R{分}{ぷん}で34.8℃、60\R{分}{ぷん}で27.3℃でした。お\R{湯}{ゆ}と\R{部屋}{へ|や}の\R{温度}{おん|ど}の\R{差}{さ}が\R{小}{ちい}さくなるほど、\R{冷}{さ}め\R{方}{かた}はゆっくりになり、\R{部屋}{へ|や}の\R{温度}{おん|ど}にだんだん\R{近}{ちか}づいていくのです（このきまりは 3-11 で\R{学}{まな}びます）。

はかった\R{範囲}{はん|い}の\R{外}{そと}へ、\R{線}{せん}をのばして\R{値}{あたい}を\R{予想}{よ|そう}することを\textbf{\R{外挿}{がい|そう}}といいます。\bkidx{がいそう}{\R{外挿}{がい|そう}}\R{外挿}{がい|そう}は、はかった\R{範囲}{はん|い}の\R{外}{そと}でもきまりが\R{変}{か}わらないときにしか、\R{正}{ただ}しくなりません。はかっていない\R{所}{ところ}では、きまりが\R{変}{か}わることがあるからです。ばねも、\R{強}{つよ}く\R{引}{ひ}きすぎると\R{比例}{ひ|れい}しなくなりました（2-15）。

\textbf{もう1つの\R{例}{れい}：\R{身長}{しん|ちょう}}\quad ある\R{子}{こ}の\R{身長}{しん|ちょう}が、10\R{歳}{さい}で140\,cm、11\R{歳}{さい}で146\,cm、12\R{歳}{さい}で152\,cmでした。1\R{年}{ねん}に6\,cmずつのびています。このまま40\R{歳}{さい}までのばすと、$140+6\times30=140+180=320$ で、320\,cm になってしまいます。\R{身長}{しん|ちょう}は、\R{大人}{おとな}になるとのびなくなるので、こうはなりません。

\begin{seikaku}
\R{外挿}{がい|そう}がまったく\R{使}{つか}えないわけではありません。\R{天気予報}{てん|き|よ|ほう}や、ロケットの\R{動}{うご}きの\R{計算}{けい|さん}は、\R{先}{さき}の\R{値}{あたい}を\R{予想}{よ|そう}する\R{外挿}{がい|そう}の\R{一種}{いっ|しゅ}です。そのときは、ただ\R{線}{せん}をのばすのではなく、\R{物}{もの}のしくみ（\R{物理}{ぶつ|り}のきまり）から\R{作}{つく}った\R{式}{しき}を\R{使}{つか}います。しくみがわかっている\R{所}{ところ}だけ、\R{先}{さき}へのばしてよいのです。
\end{seikaku}

\bksubsection{5}{グラフの\R{形}{かたち}を\R{読}{よ}む}

\leavevmode\label{bk:3-4:katachi}グラフは、\R{形}{かたち}を\R{見}{み}るだけで、\R{多}{おお}くのことを\R{教}{おし}えてくれます。ある\R{春}{はる}の\R{日}{ひ}の\R{気温}{き|おん}を、2\R{時間}{じ|かん}ごとにはかってグラフにしました。

\begin{bkfigure}
\begin{tikzpicture}
\begin{axis}[bkplot, axis lines=left, width=100mm, height=58mm, xmin=4, xmax=23.5, ymin=12, ymax=29, xtick={6,8,...,22}, ytick={14,16,...,28}, xlabel={\R{時刻}{じ|こく}（\R{時}{じ}）}, ylabel={\R{気温}{き|おん}（℃）}]
  \addplot[bkBlue, line width=1.3pt, smooth, tension=0.6] coordinates {(6,15) (8,17) (10,21) (12,24) (14,26) (16,25) (18,22) (20,20) (22,18)};
  \addplot[only marks, mark=*, mark size=1.5pt, bkInk] coordinates {(6,15) (8,17) (10,21) (12,24) (14,26) (16,25) (18,22) (20,20) (22,18)};
  \addplot[only marks, mark=*, mark size=2.4pt, bkRed] coordinates {(14,26)};
  \node[font=\sffamily\footnotesize, text=bkRed!80!black, anchor=south] at (axis cs:14,26.8) {いちばん\R{高}{たか}い（14\R{時}{じ}、26℃）};
  \node[font=\sffamily\footnotesize, text=bkGreen!40!black, anchor=north west] at (axis cs:9.3,17.6) {\R{右}{みぎ}\R{上}{あ}がり（ふえる）};
  \node[font=\sffamily\footnotesize, text=bkOrange!80!black, anchor=south west] at (axis cs:19.4,24.2) {\R{右}{みぎ}\R{下}{さ}がり（へる）};
  \node[fill=white, inner sep=0.3pt] at (axis cs:4,13.3) {\bkthrenami};
\end{axis}
\end{tikzpicture}
\bkfigcaption{ある\R{春}{はる}の\R{日}{ひ}の\R{気温}{き|おん}。}
\end{bkfigure}

\begin{itemize}[leftmargin=1.6\zw, itemsep=1pt, topsep=2pt, label={\color{\bkTT}●}]
\item \textbf{\R{右}{みぎ}\R{上}{あ}がり}の\R{所}{ところ}（6\R{時}{じ}〜14\R{時}{じ}）は、\R{時間}{じ|かん}がたつと\R{気温}{き|おん}がふえている。
\item \textbf{\R{右}{みぎ}\R{下}{さ}がり}の\R{所}{ところ}（14\R{時}{じ}〜22\R{時}{じ}）は、\R{気温}{き|おん}がへっている。
\item \R{右}{みぎ}\R{上}{あ}がりから\R{右}{みぎ}\R{下}{さ}がりに\R{変}{か}わる\R{所}{ところ}が、\textbf{いちばん\R{高}{たか}い\R{所}{ところ}}。この\R{日}{ひ}の\R{最高気温}{さい|こう|き|おん}は14\R{時}{じ}の\R{約}{やく}\qty{26}{\degreeCelsius}。
\item \R{線}{せん}が\R{急}{きゅう}な\R{所}{ところ}ほど、\R{気温}{き|おん}が\R{速}{はや}く\R{変}{か}わっている。8\R{時}{じ}〜10\R{時}{じ}は2\R{時間}{じ|かん}で4℃、14\R{時}{じ}〜16\R{時}{じ}は2\R{時間}{じ|かん}で1℃。
\end{itemize}

\textbf{2\R{本}{ほん}の\R{線}{せん}が\R{交}{まじ}わる\R{所}{ところ}}\quad \R{弟}{おとうと}が、\R{家}{いえ}から\R{図書館}{と|しょ|かん}へ、1\R{分}{ぷん}に\qty{60}{m}の\R{速}{はや}さで\R{歩}{ある}いて\R{出発}{しゅっ|ぱつ}しました。4\R{分}{ぷん}\R{後}{ご}に、\R{兄}{あに}が\R{自転車}{じ|てん|しゃ}で、1\R{分}{ぷん}に\qty{180}{m}の\R{速}{はや}さで\R{追}{お}いかけました（2-16 の\R{追}{お}いかけ）。\R{2人}{ふたり}の $x$-$t$ グラフを、\R{同}{おな}じ\R{軸}{じく}にかきます。

\noindent\begin{minipage}[c]{0.52\linewidth}\centering
\begin{tikzpicture}
\begin{axis}[bkplot, width=62mm, height=54mm, xmin=0, xmax=10.5, ymin=0, ymax=1000, xtick={0,2,...,10}, ytick={0,200,...,800}, xlabel={$t$（\R{分}{ふん}）}, ylabel={$x$（m）}]
  \addplot[bkGreen!70!black, line width=1.4pt, domain=0:10] {60*x};
  \addplot[bkBlue, line width=1.4pt, domain=4:9.1] {180*(x-4)};
  \addplot[only marks, mark=*, mark size=2.2pt, bkRed] coordinates {(6,360)};
  \draw[bkRed, line width=0.6pt, dash pattern=on 2pt off 1.5pt] (axis cs:6,0) -- (axis cs:6,360) -- (axis cs:0,360);
  \node[font=\sffamily\footnotesize, text=bkGreen!40!black, anchor=west] at (axis cs:9.2,600) {\R{弟}{おとうと}};
  \node[font=\sffamily\footnotesize, text=bkBlue!80!black, anchor=west] at (axis cs:8.2,920) {\R{兄}{あに}};
\end{axis}
\end{tikzpicture}
\end{minipage}\hfill
\begin{minipage}[c]{0.45\linewidth}
\R{弟}{おとうと}の\R{線}{せん}（\R{緑}{みどり}）は0\R{分}{ふん}から、\R{兄}{あに}の\R{線}{せん}（\R{青}{あお}）は4\R{分}{ぷん}から\R{始}{はじ}まります。\R{兄}{あに}の\R{線}{せん}のほうが\R{急}{きゅう}です。2\R{本}{ほん}の\R{線}{せん}は、6\R{分}{ぷん}・360\,mの\R{点}{てん}で\R{交}{まじ}わります。

\R{交}{まじ}わる\R{点}{てん}は「\R{同}{おな}じ\R{時刻}{じ|こく}に、\R{同}{おな}じ\R{場所}{ば|しょ}にいる」\R{点}{てん}なので、ここで\R{兄}{あに}が\R{弟}{おとうと}に\textbf{\R{追}{お}いつく}のです。
\end{minipage}

\par\vspace{3pt}
\R{確}{たし}かめてみましょう。6\R{分}{ぷん}のとき、\R{家}{いえ}からの\R{位置}{い|ち}は、
\begin{itemize}[leftmargin=1.6\zw, itemsep=0pt, topsep=1pt, label={\color{\bkTT}●}]
\item \R{弟}{おとうと}：$60\times6=360$（m）
\item \R{兄}{あに}：\R{走}{はし}ったのは $6-4=2$（\R{分}{ふん}）なので、$180\times2=360$（m）
\end{itemize}
で、2\R{人}{にん}は\R{同}{おな}じ\R{所}{ところ}にいます。\R{交}{まじ}わる\R{点}{てん}を、グラフを\R{読}{よ}まずに\R{式}{しき}で\R{求}{もと}める\R{方法}{ほう|ほう}は、3-6 で\R{学}{まな}びます。

\begin{kotoba}
\small
\bkword{$x$-$t$ グラフ}{\R{横軸}{よこ|じく}を\R{時間}{じ|かん} $t$、たて\R{軸}{じく}を\R{位置}{い|ち} $x$ にしたグラフ。\R{水平}{すい|へい}なら\R{止}{と}まっている。\R{急}{きゅう}なほど\R{速}{はや}い。}
\bkword{$v$-$t$ グラフ}{\R{横軸}{よこ|じく}を\R{時間}{じ|かん} $t$、たて\R{軸}{じく}を\R{速}{はや}さ $v$ にしたグラフ（2-3）。\R{線}{せん}の\R{下}{した}の\R{面積}{めん|せき}が\R{進}{すす}んだ\R{道}{みち}のり。}
\bkword{\R{内挿}{ない|そう}}{はかった\R{点}{てん}と\R{点}{てん}のあいだの\R{値}{あたい}を、グラフの\R{線}{せん}から\R{読}{よ}みとること。}
\bkword{\R{外挿}{がい|そう}}{はかった\R{範囲}{はん|い}の\R{外}{そと}へ\R{線}{せん}をのばして、\R{値}{あたい}を\R{予想}{よ|そう}すること。きまりが\R{変}{か}わることがあるので\R{危}{あぶ}ない。}
\end{kotoba}

%% ---------------- 例題 ----------------
\bksection{reidai}{\R{例題}{れい|だい}でたしかめよう}

\begin{reidai}[\R{軸}{じく}と\R{目}{め}もりを\R{決}{き}めてグラフをかく]
\R{水}{みず}を\R{電熱線}{でん|ねつ|せん}であたためて、1\R{分}{ぷん}ごとに\R{温度}{おん|ど}をはかりました。\R{方眼}{ほう|がん}の\R{横}{よこ}は12マス、たては10マスです。\R{横軸}{よこ|じく}とたて\R{軸}{じく}に\R{何}{なに}をとり、1マスをいくつにすればよいですか。グラフもかきましょう（ノートに）。\par
\begin{bkcenter}
\begin{bktable}{colspec={Q[l,m,wd=20mm] *{7}{Q[c,m,wd=9mm]}}, rowsep=1pt, cells={font=\small}}
\R{時間}{じ|かん}（\R{分}{ふん}） & 0 & 1 & 2 & 3 & 4 & 5 & 6 \\
\R{温度}{おん|ど}（℃） & 20.0 & 22.7 & 25.2 & 27.9 & 30.4 & 33.1 & 35.6 \\
\end{bktable}
\end{bkcenter}
\kangaekata \R{時間}{じ|かん}を\R{横軸}{よこ|じく}、\R{温度}{おん|ど}をたて\R{軸}{じく}にとります。\R{全部}{ぜん|ぶ}の\R{値}{あたい}が\R{入}{はい}り、できるだけ\R{大}{おお}きくかける\R{目}{め}もりを、1・2・5…から\R{選}{えら}びます。\par
\bkthreSi（1）\nobreak \R{横軸}{よこ|じく}：\R{時間}{じ|かん} $t$（\R{分}{ふん}）。0〜6\R{分}{ぷん}を12マスに\R{入}{い}れるので、1マス $6\div12=0.5$（\R{分}{ふん}）
\bkthreI（2）\nobreak たて\R{軸}{じく}：\R{温度}{おん|ど} $T$（℃）。20.0〜35.6℃を10マスに\R{入}{い}れる。
\bkthreJ 0から\R{始}{はじ}めると、1マス4℃は\R{読}{よ}みにくいので5℃。0〜50℃で、\R{点}{てん}は\R{下}{した}\R{半分}{はん|ぶん}にかたまる。
\bkthreJ 20℃より\R{下}{した}を \bkthrenami で\R{省}{はぶ}くと、1マス2℃で20〜40℃。\R{点}{てん}が\R{全体}{ぜん|たい}に\R{広}{ひろ}がる。
\bkthreI（3）\nobreak \R{点}{てん}はほぼまっすぐに\R{並}{なら}ぶので、\R{定規}{じょう|ぎ}で\R{直線}{ちょく|せん}を\R{引}{ひ}く（グラフ）。\par
\noindent\begin{minipage}[c]{0.5\linewidth}
\bkthreK \R{横軸}{よこ|じく}は\R{時間}{じ|かん}で1マス0.5\R{分}{ふん}。たて\R{軸}{じく}は\R{温度}{おん|ど}で、20℃より\R{下}{した}を\R{省}{はぶ}いて1マス2℃。
\end{minipage}\hfill
\begin{minipage}[c]{0.46\linewidth}\centering
\begin{tikzpicture}
\begin{axis}[bkplot, axis lines=left, width=46mm, height=44mm, xmin=0, xmax=6, ymin=20, ymax=40, xtick={0,1,...,6}, ytick={20,24,...,40}, minor xtick={0.5,1.5,...,5.5}, minor ytick={22,26,...,38}, grid=both, minor grid style={bkLine, line width=0.3pt}, xlabel={$t$（\R{分}{ふん}）}, ylabel={$T$（℃）}]
  \addplot[bkBlue, line width=1.2pt, domain=0:6] {20+2.6*x};
  \addplot[only marks, mark=*, mark size=1.4pt, bkInk] coordinates {(0,20.0) (1,22.7) (2,25.2) (3,27.9) (4,30.4) (5,33.1) (6,35.6)};
\end{axis}
\end{tikzpicture}
\end{minipage}
\end{reidai}


\begin{reidai}[$x$-$t$ グラフを\R{読}{よ}む]
\noindent\begin{minipage}[c]{0.55\linewidth}
ミナさんが\R{家}{いえ}から\R{公園}{こう|えん}へ\R{行}{い}き、しばらく\R{遊}{あそ}んでから\R{家}{いえ}へ\R{帰}{かえ}りました。グラフは、そのときの $x$-$t$ グラフです。（1）\R{公園}{こう|えん}は\R{家}{いえ}から\R{何}{なん}mですか。（2）\R{公園}{こう|えん}には\R{何分}{なん|ぷん}いましたか。（3）\R{行}{い}きと\R{帰}{かえ}りの\R{速}{はや}さは、それぞれ1\R{分}{ぷん}に\R{何}{なん}mですか。（4）\R{帰}{かえ}りの\R{線}{せん}が\R{右}{みぎ}\R{下}{さ}がりなのは、なぜですか。
\end{minipage}\hfill
\begin{minipage}[c]{0.42\linewidth}\centering
\begin{tikzpicture}
\begin{axis}[bkplot, width=42mm, height=38mm, xmin=0, xmax=27, ymin=0, ymax=640, xtick={0,5,...,25}, ytick={0,100,...,500}, xlabel={$t$（\R{分}{ふん}）}, ylabel={$x$（m）}]
  \addplot[bkBlue, line width=1.3pt] coordinates {(0,0) (10,500) (20,500) (25,0)};
\end{axis}
\end{tikzpicture}
\end{minipage}\par
\kangaekata \R{線}{せん}の\R{高}{たか}さが\R{位置}{い|ち}、\R{水平}{すい|へい}な\R{所}{ところ}が\R{止}{と}まっている\R{所}{ところ}です。\R{速}{はや}さは、\R{進}{すす}んだ\R{道}{みち}のり÷かかった\R{時間}{じ|かん}（2-3）で\R{求}{もと}めます。\par
\bkthreSi（1）\nobreak \R{線}{せん}がいちばん\R{高}{たか}い\R{所}{ところ}の\R{高}{たか}さは500。
\bkthreI（2）\nobreak \R{線}{せん}が\R{水平}{すい|へい}なのは10\R{分}{ぷん}から20\R{分}{ぷん}。$20-10=10$（\R{分}{ふん}）
\bkthreI（3）\nobreak \R{行}{い}き：10\R{分}{ぷん}で500\,m\R{進}{すす}んだので、$500\div10=50$
\bkthreJ \R{帰}{かえ}り：20\R{分}{ぷん}から25\R{分}{ふん}までの5\R{分}{ふん}で500\,m\R{進}{すす}んだので、$500\div5=100$
\bkthreI（4）\nobreak \R{家}{いえ}に\R{近}{ちか}づくと、\R{家}{いえ}からの\R{位置}{い|ち} $x$ が\R{小}{ちい}さくなるから。\par
\bkthreK（1）\nobreak \qty{500}{m}\quad（2）\nobreak 10\R{分}{ぷん}\quad（3）\nobreak \R{行}{い}き\qty{50}{m}、\R{帰}{かえ}り\qty{100}{m}\quad（4）\nobreak \R{位置}{い|ち} $x$ がへっていくから
\end{reidai}

\begingroup\tcbset{equal height group=bkthrejA}%
\noindent\begin{minipage}[t]{0.49\linewidth}
\begin{jibun}
ろうそくをもやして、2\R{分}{ふん}ごとに\R{長}{なが}さをはかりました。\R{方眼}{ほう|がん}は\R{横}{よこ}10マス、たて10マスです。\R{横軸}{よこ|じく}とたて\R{軸}{じく}に\R{何}{なに}をとり、1マスをいくつにすればよいですか。グラフもかきましょう（ノートに）。\par\smallskip
{\small \R{時間}{じ|かん}（\R{分}{ふん}）0、2、4、6、8、10\par \R{長}{なが}さ（cm）15.0、13.9、13.1、12.0、10.9、10.1}
\bkjans{\R{横軸}{よこ|じく}は\R{時間}{じ|かん} $t$（\R{分}{ふん}）で1マス1\R{分}{ぷん}。たて\R{軸}{じく}は\R{長}{なが}さ（cm）。0からだと1マス2\,cm、10\,cmより\R{下}{した}を \bkthrenami で\R{省}{はぶ}くと1マス0.5\,cmで10〜15\,cm（こちらが\R{大}{おお}きくかける）。\R{点}{てん}がまっすぐに\R{並}{なら}ぶので、\R{右}{みぎ}\R{下}{さ}がりの\R{直線}{ちょく|せん}を\R{引}{ひ}く}
\end{jibun}
\end{minipage}\hfill
\begin{minipage}[t]{0.49\linewidth}
\begin{jibun}
ソラさんが\R{家}{いえ}からコンビニへ\R{行}{い}って、\R{家}{いえ}へ\R{帰}{かえ}りました。（1）コンビニは\R{家}{いえ}から\R{何}{なん}mですか。（2）コンビニに\R{何分}{なん|ぷん}いましたか。（3）\R{行}{い}きと\R{帰}{かえ}りの\R{速}{はや}さは、1\R{分}{ぷん}に\R{何}{なん}mですか。\par\smallskip
{\centering
\begin{tikzpicture}
\begin{axis}[bkplot, width=44mm, height=30mm, xmin=0, xmax=22, ymin=0, ymax=760, xtick={0,4,...,20}, ytick={0,200,400,600}, xlabel={$t$（\R{分}{ふん}）}, ylabel={$x$（m）}]
  \addplot[bkBlue, line width=1.3pt] coordinates {(0,0) (8,600) (14,600) (20,0)};
\end{axis}
\end{tikzpicture}\par}
\bkjans{（1）\nobreak \qty{600}{m}\quad（2）\nobreak $14-8=6$ で6\R{分}{ぷん}\quad（3）\nobreak \R{行}{い}き $600\div8=75$ で\qty{75}{m}、\R{帰}{かえ}り $600\div6=100$ で\qty{100}{m}}
\end{jibun}
\end{minipage}\par
\endgroup

\begin{column}{グラフをかいたら、5つのチェック}
\begin{enumerate}[label={\bkcircled[bkNavy]{\arabic*}}, leftmargin=2.4\zw, itemsep=1pt, topsep=1pt]
\item \R{横軸}{よこ|じく}とたて\R{軸}{じく}に、\R{量}{りょう}の\R{名前}{な|まえ}と\R{単位}{たん|い}を\R{書}{か}いたかな？（「$t$（\R{分}{ふん}）」「$T$（℃）」）
\item 1\R{目}{め}もりは、1・2・5（10・20・50、0.1・0.2・0.5）のどれかかな？
\item \R{目}{め}もりは、\R{等}{ひと}しい\R{間}{ま}かくでならんでいるかな？
\item 0の\R{近}{ちか}くを\R{省}{はぶ}いたときは、\R{軸}{じく}に \bkthrenami をかいたかな？
\item \R{実験}{じっ|けん}の\R{点}{てん}には、おれ\R{線}{せん}ではなく、なめらかな\R{線}{せん}を1\R{本}{ぽん}\R{引}{ひ}いたかな？
\end{enumerate}
\end{column}



\begin{reidai}[$v$-$t$ グラフから\R{位置}{い|ち}を\R{求}{もと}める]
\noindent\begin{minipage}[c]{0.55\linewidth}
ロボットが、まっすぐな\R{線}{せん}の\R{上}{うえ}を\R{動}{うご}きました。グラフは $v$-$t$ グラフです。\R{動}{うご}きはじめの\R{位置}{い|ち}を0とします。（1）4\R{秒}{びょう}、6\R{秒}{びょう}、10\R{秒}{びょう}のときの\R{位置}{い|ち}を\R{求}{もと}めましょう。（2）$x$-$t$ グラフの\R{形}{かたち}を\R{言}{い}いましょう。
\end{minipage}\hfill
\begin{minipage}[c]{0.42\linewidth}\centering
\begin{tikzpicture}
\begin{axis}[bkplot, width=46mm, height=32mm, xmin=0, xmax=11, ymin=0, ymax=6.5, xtick={0,2,...,10}, ytick={0,3,5}, xlabel={$t$（s）}, ylabel={$v$（m/s）}]
  \addplot[fill=bkBlue!15, draw=none] coordinates {(0,0) (0,3) (4,3) (4,0)} \closedcycle;
  \addplot[fill=bkBlue!15, draw=none] coordinates {(6,0) (6,5) (10,5) (10,0)} \closedcycle;
  \addplot[bkOrange, line width=1.3pt] coordinates {(0,3) (4,3)};
  \addplot[bkOrange, line width=1.3pt] coordinates {(4,0) (6,0)};
  \addplot[bkOrange, line width=1.3pt] coordinates {(6,5) (10,5)};
\end{axis}
\end{tikzpicture}
\end{minipage}\par
\kangaekata \R{進}{すす}んだ\R{道}{みち}のりは、$v$-$t$ グラフの\R{線}{せん}の\R{下}{した}の\R{長方形}{ちょう|ほう|けい}の\R{面積}{めん|せき}（\R{速}{はや}さ×\R{時間}{じ|かん}）です（2-3）。\R{位置}{い|ち}は、\R{前}{まえ}の\R{位置}{い|ち}に、\R{進}{すす}んだ\R{道}{みち}のりをたしていきます。\par
\bkthreSi（1）\nobreak 0〜4\R{秒}{びょう}：$3\times4=12$。4\R{秒}{びょう}の\R{位置}{い|ち}は $0+12=12$
\bkthreJ 4〜6\R{秒}{びょう}：\R{速}{はや}さ0で\R{止}{と}まっている。6\R{秒}{びょう}の\R{位置}{い|ち}は12のまま
\bkthreJ 6〜10\R{秒}{びょう}：$5\times4=20$。10\R{秒}{びょう}の\R{位置}{い|ち}は $12+20=32$
\bkthreI（2）\nobreak 0〜4\R{秒}{びょう}は\R{右}{みぎ}\R{上}{あ}がりの\R{直線}{ちょく|せん}、4〜6\R{秒}{びょう}は\R{水平}{すい|へい}、6〜10\R{秒}{びょう}は\R{前}{まえ}より\R{急}{きゅう}な\R{右}{みぎ}\R{上}{あ}がりの\R{直線}{ちょく|せん}。\par
\bkthreK（1）\nobreak 4\R{秒}{びょう}で\qty{12}{m}、6\R{秒}{びょう}で\qty{12}{m}、10\R{秒}{びょう}で\qty{32}{m}\quad（2）\nobreak \R{上}{あ}がる→\R{水平}{すい|へい}→もっと\R{急}{きゅう}に\R{上}{あ}がる
\end{reidai}

\begin{jibun}
\noindent\begin{minipage}[c]{0.5\linewidth}
\R{台車}{だい|しゃ}の $v$-$t$ グラフです。\R{動}{うご}きはじめの\R{位置}{い|ち}を0とします。（1）5\R{秒}{びょう}、7\R{秒}{びょう}、10\R{秒}{びょう}のときの\R{位置}{い|ち}を\R{求}{もと}めましょう。（2）$x$-$t$ グラフの\R{形}{かたち}を\R{言}{い}いましょう。
\end{minipage}\hfill
\begin{minipage}[c]{0.46\linewidth}\centering
\begin{tikzpicture}
\begin{axis}[bkplot, width=52mm, height=34mm, xmin=0, xmax=11, ymin=0, ymax=5.4, xtick={0,2,...,10}, ytick={0,2,4}, xlabel={$t$（s）}, ylabel={$v$（m/s）}]
  \addplot[bkOrange, line width=1.3pt] coordinates {(0,2) (5,2)};
  \addplot[bkOrange, line width=1.3pt] coordinates {(5,0) (7,0)};
  \addplot[bkOrange, line width=1.3pt] coordinates {(7,4) (10,4)};
\end{axis}
\end{tikzpicture}
\end{minipage}
\bkjans{（1）\nobreak 5\R{秒}{びょう}は $2\times5=10$ で\qty{10}{m}、7\R{秒}{びょう}も\qty{10}{m}、10\R{秒}{びょう}は $10+4\times3=22$ で\qty{22}{m}\quad（2）\nobreak \R{右}{みぎ}\R{上}{あ}がりの\R{直線}{ちょく|せん}→\R{水平}{すい|へい}→もっと\R{急}{きゅう}な\R{右}{みぎ}\R{上}{あ}がりの\R{直線}{ちょく|せん}}
\end{jibun}

\begin{reidai}[\R{内挿}{ない|そう}と\R{外挿}{がい|そう}]
やかんの\R{水}{みず}を\R{熱}{ねっ}して、2\R{分}{ふん}ごとに\R{温度}{おん|ど}をはかりました。\R{点}{てん}はまっすぐに\R{並}{なら}びました。\par
\begin{bkcenter}
\begin{bktable}{colspec={Q[l,m,wd=20mm] *{5}{Q[c,m,wd=10mm]}}, rowsep=1pt, cells={font=\small}}
\R{時間}{じ|かん}（\R{分}{ふん}） & 0 & 2 & 4 & 6 & 8 \\
\R{温度}{おん|ど}（℃） & 20 & 36 & 52 & 68 & 84 \\
\end{bktable}
\end{bkcenter}
（1）5\R{分}{ふん}のときの\R{温度}{おん|ど}を\R{求}{もと}めましょう。（2）76℃になったのは\R{何分}{なん|ぷん}のときですか。（3）この\R{直線}{ちょく|せん}を12\R{分}{ふん}までのばすと\R{何}{なん}℃になりますか。それは\R{正}{ただ}しいと\R{言}{い}えますか。\par
\kangaekata 2\R{分}{ふん}ごとに16℃ずつ\R{上}{あ}がっているので、1\R{分}{ぷん}に $16\div2=8$（℃）\R{上}{あ}がります。（1）（2）ははかった\R{範囲}{はん|い}の\R{中}{なか}（\R{内挿}{ない|そう}）、（3）は\R{外}{そと}（\R{外挿}{がい|そう}）です。\par
\bkthreSi（1）\nobreak 5\R{分}{ふん}は4\R{分}{ぷん}と6\R{分}{ぷん}のまん\R{中}{なか}。$(52+68)\div2=120\div2=60$
\bkthreI（2）\nobreak 76は68と84のまん\R{中}{なか}なので、6\R{分}{ぷん}と8\R{分}{ぷん}のまん\R{中}{なか}の7\R{分}{ふん}
\bkthreI（3）\nobreak 0\R{分}{ふん}に20℃だった\R{水}{みず}が、1\R{分}{ぷん}に8℃ずつ12\R{分}{ふん}\R{上}{あ}がると、$20+8\times12=20+96=116$
\bkthreJ でも\R{水}{みず}は\qty{100}{\degreeCelsius}でふっとうし、それより\R{熱}{あつ}くならない。\par
\bkthreK（1）\nobreak \qty{60}{\degreeCelsius}\quad（2）\nobreak 7\R{分}{ふん}\quad（3）\nobreak \qty{116}{\degreeCelsius}になるが、\R{正}{ただ}しくない（\R{水}{みず}は\qty{100}{\degreeCelsius}より\R{熱}{あつ}くならない。はかった\R{範囲}{はん|い}の\R{外}{そと}では、きまりが\R{変}{か}わる）
\end{reidai}

\begin{jibun}
ある\R{子}{こ}の\R{身長}{しん|ちょう}をはかると、8\R{歳}{さい}で126\,cm、9\R{歳}{さい}で132\,cm、10\R{歳}{さい}で138\,cm、11\R{歳}{さい}で144\,cmで、\R{点}{てん}はまっすぐに\R{並}{なら}びました。（1）9\R{歳}{さい}\R{半}{はん}のときの\R{身長}{しん|ちょう}を\R{求}{もと}めましょう。（2）141\,cmだったのは\R{何歳}{なん|さい}のときですか。（3）この\R{直線}{ちょく|せん}を30\R{歳}{さい}までのばすと\R{何}{なん}cmになりますか。それは\R{正}{ただ}しいと\R{言}{い}えますか。
\bkjans{（1）\nobreak $(132+138)\div2=135$ で\qty{135}{cm}\quad（2）\nobreak 141は138と144のまん\R{中}{なか}なので10\R{歳}{さい}\R{半}{はん}\quad（3）\nobreak 1\R{年}{ねん}に6\,cmで22\R{年}{ねん}\R{分}{ぶん}、$126+6\times22=258$ で\qty{258}{cm}。\R{正}{ただ}しくない（\R{大人}{おとな}になると\R{身長}{しん|ちょう}はのびなくなる）}
\end{jibun}

\begin{reidai}[2\R{本}{ほん}の\R{線}{せん}が\R{交}{まじ}わる\R{所}{ところ}を\R{読}{よ}む]
\R{弟}{おとうと}が\R{家}{いえ}から\R{図書館}{と|しょ|かん}へ、1\R{分}{ぷん}に\qty{60}{m}で\R{歩}{ある}いて\R{出発}{しゅっ|ぱつ}しました。4\R{分}{ぷん}\R{後}{ご}に、\R{兄}{あに}が1\R{分}{ぷん}に\qty{180}{m}の\R{自転車}{じ|てん|しゃ}で\R{追}{お}いかけました（\pageref{bk:3-4:katachi}ページのグラフ）。\R{図書館}{と|しょ|かん}は\R{家}{いえ}から\qty{900}{m}です。（1）\R{兄}{あに}が\R{出発}{しゅっ|ぱつ}したとき、\R{弟}{おとうと}は\R{家}{いえ}から\R{何}{なん}mの\R{所}{ところ}にいましたか。（2）\R{兄}{あに}が\R{弟}{おとうと}に\R{追}{お}いつくのは、\R{何分}{なん|ぷん}の、\R{家}{いえ}から\R{何}{なん}mの\R{所}{ところ}ですか。（3）どちらが\R{何分}{なん|ぷん}\R{早}{はや}く\R{図書館}{と|しょ|かん}に\R{着}{つ}きますか。\par
\kangaekata（2）は2\R{本}{ほん}の\R{線}{せん}の\R{交}{まじ}わる\R{点}{てん}を\R{読}{よ}み、\R{計算}{けい|さん}で\R{確}{たし}かめます。（3）は、かかる\R{時間}{じ|かん}＝\R{道}{みち}のり÷\R{速}{はや}さ（2-14）で\R{求}{もと}めます。\par
\bkthreSi（1）\nobreak $60\times4=240$
\bkthreI（2）\nobreak \R{交}{まじ}わる\R{点}{てん}は6\R{分}{ぷん}・360\,m。\R{確}{たし}かめ：\R{弟}{おとうと} $60\times6=360$、\R{兄}{あに} $180\times2=360$
\bkthreI（3）\nobreak \R{弟}{おとうと}：$900\div60=15$（\R{分}{ふん}）
\bkthreJ \R{兄}{あに}：$900\div180=5$。4\R{分}{ぷん}\R{後}{ご}に\R{出}{で}たので $4+5=9$（\R{分}{ふん}）
\bkthreJ $15-9=6$\par
\bkthreK（1）\nobreak \qty{240}{m}\quad（2）\nobreak 6\R{分}{ぷん}、\R{家}{いえ}から\qty{360}{m}\quad（3）\nobreak \R{兄}{あに}が6\R{分}{ぷん}\R{早}{はや}い
\end{reidai}

\begin{jibun}
Aさんは\R{家}{いえ}から1\R{分}{ぷん}に\qty{70}{m}で\R{歩}{ある}いて\R{出発}{しゅっ|ぱつ}し、6\R{分}{ぷん}\R{後}{ご}にBさんが1\R{分}{ぷん}に\qty{210}{m}の\R{自転車}{じ|てん|しゃ}で\R{追}{お}いかけました。\R{駅}{えき}は\R{家}{いえ}から\qty{1050}{m}です。（1）Bさんが\R{出発}{しゅっ|ぱつ}したとき、Aさんは\R{家}{いえ}から\R{何}{なん}mですか。（2）Bさんが\R{追}{お}いつくのは\R{何分}{なん|ぷん}の、\R{家}{いえ}から\R{何}{なん}mの\R{所}{ところ}ですか（ノートにグラフをかいて\R{読}{よ}み、\R{計算}{けい|さん}で\R{確}{たし}かめよう）。（3）どちらが\R{何分}{なん|ぷん}\R{早}{はや}く\R{駅}{えき}に\R{着}{つ}きますか。
\bkjans{（1）\nobreak $70\times6=420$ で\qty{420}{m}\quad（2）\nobreak 9\R{分}{ふん}、\qty{630}{m}（Aさん $70\times9=630$、Bさん $210\times3=630$）\quad（3）\nobreak Aさん $1050\div70=15$ \R{分}{ふん}、Bさん $6+1050\div210=6+5=11$ \R{分}{ふん}。Bさんが4\R{分}{ぷん}\R{早}{はや}い}
\end{jibun}

\begin{reidai}[1\R{目}{め}もりのちがうグラフをくらべる]
AさんとBさんの $x$-$t$ グラフです。Bさんの\R{線}{せん}のほうが\R{急}{きゅう}に\R{見}{み}えます。どちらが\R{速}{はや}いですか。\par
\noindent\begin{minipage}[c]{0.49\linewidth}\centering
\begin{tikzpicture}
\begin{axis}[bkplot, width=52mm, height=44mm, xmin=0, xmax=11, ymin=0, ymax=86, xtick={0,2,...,10}, ytick={0,10,...,80}, xlabel={$t$（s）}, ylabel={$x$（m）}]
  \addplot[bkBlue, line width=1.3pt] coordinates {(0,0) (10,50)};
  \node[font=\sffamily\bfseries\footnotesize, anchor=south] at (axis cs:5,75) {Aさん};
\end{axis}
\end{tikzpicture}
\end{minipage}\hfill
\begin{minipage}[c]{0.49\linewidth}\centering
\begin{tikzpicture}
\begin{axis}[bkplot, width=52mm, height=44mm, xmin=0, xmax=11, ymin=0, ymax=43, xtick={0,2,...,10}, ytick={0,5,...,40}, xlabel={$t$（s）}, ylabel={$x$（m）}]
  \addplot[bkRed, line width=1.3pt] coordinates {(0,0) (10,40)};
  \node[font=\sffamily\bfseries\footnotesize, anchor=south] at (axis cs:3,34) {Bさん};
\end{axis}
\end{tikzpicture}
\end{minipage}\par
\kangaekata \R{見}{み}た\R{目}{め}の\R{急}{きゅう}さは、1\R{目}{め}もりの\R{大}{おお}きさで\R{変}{か}わります。\R{目}{め}もりの\R{数字}{すう|じ}を\R{読}{よ}んで、\R{速}{はや}さ＝\R{進}{すす}んだ\R{道}{みち}のり÷\R{時間}{じ|かん}を\R{計算}{けい|さん}します。\par
\bkthreS Aさん：たて1\R{目}{め}もり10\,m。10\R{秒}{びょう}で50\,m。$50\div10=5$
\bkthreC Bさん：たて1\R{目}{め}もり5\,m。10\R{秒}{びょう}で40\,m。$40\div10=4$\par
\bkthreK Aさんが\R{速}{はや}い（Aさんは\qty{5}{m/s}、Bさんは\qty{4}{m/s}）
\end{reidai}

\begin{jibun}
CさんとDさんの $x$-$t$ グラフです。Dさんの\R{線}{せん}のほうが\R{急}{きゅう}に\R{見}{み}えます。どちらが\R{速}{はや}いですか。\par\smallskip
\centering
\begin{tikzpicture}
\begin{axis}[bkplot, name=c, width=40mm, height=44mm, xmin=0, xmax=66, ymin=0, ymax=660, xtick={0,30,60}, ytick={0,100,...,600}, xlabel={$t$（s）}, ylabel={$x$（m）}, every axis y label/.append style={font=\sffamily\scriptsize}, every axis x label/.append style={font=\sffamily\scriptsize}]
  \addplot[bkBlue, line width=1.2pt] coordinates {(0,0) (60,300)};
  \node[font=\sffamily\bfseries\scriptsize, anchor=north west] at (axis cs:2,650) {C};
\end{axis}
\begin{axis}[bkplot, at={($(c.east)+(16mm,0)$)}, anchor=west, width=40mm, height=44mm, xmin=0, xmax=66, ymin=0, ymax=198, xtick={0,30,60}, ytick={0,30,...,180}, xlabel={$t$（s）}, ylabel={$x$（m）}, every axis y label/.append style={font=\sffamily\scriptsize}, every axis x label/.append style={font=\sffamily\scriptsize}]
  \addplot[bkRed, line width=1.2pt] coordinates {(0,0) (60,180)};
  \node[font=\sffamily\bfseries\scriptsize, anchor=north west] at (axis cs:2,195) {D};
\end{axis}
\end{tikzpicture}
\bkjans{Cさんが\R{速}{はや}い（Cさんは $300\div60=5$ で\qty{5}{m/s}、Dさんは $180\div60=3$ で\qty{3}{m/s}）}
\end{jibun}

%% ---------------- まちがえやすいところ・物理工学 ----------------
\bksection{tsumazuki}{まちがえやすいところをチェック}\label{bk:3-4:tsuma}

\begin{tsumazuki}{1\R{目}{め}もりのちがうグラフを、\R{見}{み}た\R{目}{め}でくらべる}
\bkng{\R{線}{せん}が\R{急}{きゅう}に\R{見}{み}えるほうが、いつも\R{速}{はや}い}
\bkok{\R{目}{め}もりの\R{数字}{すう|じ}を\R{読}{よ}んで、\R{道}{みち}のり÷\R{時間}{じ|かん}を\R{計算}{けい|さん}してくらべる}
たて\R{軸}{じく}の1\R{目}{め}もりを\R{半分}{はん|ぶん}にすると、\R{同}{おな}じ\R{速}{はや}さでも\R{線}{せん}は2\R{倍}{ばい}\R{急}{きゅう}に\R{見}{み}えます（1-18）。
\end{tsumazuki}

\begin{tsumazuki}{$x$-$t$ グラフの\R{水平}{すい|へい}な\R{線}{せん}}
\bkng{$x$-$t$ グラフの\R{水平}{すい|へい}な\R{線}{せん}は「\R{同}{おな}じ\R{速}{はや}さで\R{進}{すす}んでいる」}
\bkok{$x$-$t$ グラフの\R{水平}{すい|へい}な\R{線}{せん}は「\R{止}{と}まっている」}
\R{位置}{い|ち}が\R{変}{か}わらないから\R{水平}{すい|へい}です。$v$-$t$ グラフの\R{水平}{すい|へい}な\R{線}{せん}（\R{速}{はや}さが\R{一定}{いっ|てい}）と\R{混}{ま}ぜないようにしよう。
\end{tsumazuki}

\begin{tsumazuki}{はかった\R{時刻}{じ|こく}を、\R{等}{ひと}しい\R{間}{かん}かくに\R{並}{なら}べる}
\bkng{0\R{分}{ふん}・1\R{分}{ぷん}・2\R{分}{ふん}・5\R{分}{ふん}・10\R{分}{ぷん}を、\R{横軸}{よこ|じく}に\R{同}{おな}じ\R{間}{かん}かくでならべる}
\bkok{\R{横軸}{よこ|じく}は1\R{目}{め}もりを\R{決}{き}めて、5\R{分}{ふん}は5の\R{所}{ところ}、10\R{分}{ぷん}は10の\R{所}{ところ}にかく}
\R{目}{め}もりが\R{等}{ひと}しい\R{間}{かん}かくでないと、グラフの\R{形}{かたち}がゆがみ、\R{変}{か}わり\R{方}{かた}を\R{読}{よ}みまちがえます。
\end{tsumazuki}

\begin{tsumazuki}{\R{実験}{じっ|けん}の\R{点}{てん}をおれ\R{線}{せん}でつなぐ}
\bkng{はかった\R{点}{てん}を、\R{全部}{ぜん|ぶ}\R{通}{とお}るおれ\R{線}{せん}で\R{結}{むす}ぶ}
\bkok{\R{点}{てん}の\R{近}{ちか}くを\R{通}{とお}る、なめらかな\R{線}{せん}を1\R{本}{ぽん}\R{引}{ひ}く}
\R{点}{てん}のでこぼこは、はかるときのずれです。\R{全部}{ぜん|ぶ}の\R{点}{てん}を\R{通}{とお}さなくてよいのです。
\end{tsumazuki}

\begin{tsumazuki}{\R{線}{せん}をどこまでものばす}
\bkng{はかった\R{範囲}{はん|い}の\R{外}{そと}も、\R{線}{せん}をのばした\R{値}{あたい}になる}
\bkok{\R{範囲}{はん|い}の\R{外}{そと}では、きまりが\R{変}{か}わることがある}
お\R{湯}{ゆ}は\R{部屋}{へ|や}の\R{温度}{おん|ど}より\R{冷}{つめ}たくならず、\R{水}{みず}は\qty{100}{\degreeCelsius}より\R{熱}{あつ}くならず、\R{身長}{しん|ちょう}はのびなくなります。
\end{tsumazuki}

\begin{column}{\R{電車}{でん|しゃ}のダイヤグラムは $x$-$t$ グラフ}
\noindent\begin{minipage}[c]{0.58\linewidth}
\R{電車}{でん|しゃ}の\R{会社}{かい|しゃ}は、\R{横軸}{よこ|じく}を\R{時刻}{じ|こく}、たて\R{軸}{じく}を\R{駅}{えき}の\R{位置}{い|ち}にした $x$-$t$ グラフで、\R{電車}{でん|しゃ}の\R{動}{うご}きを\R{計画}{けい|かく}します。このグラフを\textbf{ダイヤグラム}といいます。\bkidx{だいやぐらむ}{ダイヤグラム}\R{水平}{すい|へい}な\R{所}{ところ}は\R{駅}{えき}に\R{止}{と}まっている\R{所}{ところ}、\R{急}{きゅう}な\R{線}{せん}は\R{速}{はや}い\R{特急}{とっ|きゅう}、\R{交}{まじ}わる\R{所}{ところ}は2\R{本}{ほん}の\R{電車}{でん|しゃ}がすれちがったり\R{追}{お}いこしたりする\R{所}{ところ}です。「ダイヤがみだれる」の「ダイヤ」は、このグラフのことです。
\end{minipage}\hfill
\begin{minipage}[c]{0.4\linewidth}\centering
\begin{tikzpicture}[x=0.82mm, y=1mm]
  \foreach \y/\s in {0/A\R{駅}{えき},10/B\R{駅}{えき},20/C\R{駅}{えき},30/D\R{駅}{えき}} {\draw[bkLine, line width=0.4pt] (10,\y) -- (58,\y); \node[font=\sffamily\scriptsize, anchor=east] at (9,\y) {\s};}
  \foreach \k in {0,...,6} {\draw[bkLine, line width=0.3pt] (10+\k*8,0) -- (10+\k*8,30);}
  \draw[-{Stealth[length=3pt]}, bkInk, line width=0.6pt] (10,0) -- (60,0);
  \node[font=\sffamily\scriptsize, anchor=north] at (56,-0.5) {\R{時刻}{じ|こく}};
  \draw[bkBlue, line width=1.1pt] (10,0) -- (19,10) -- (22,10) -- (31,20) -- (36,20) -- (45,30);
  \draw[bkRed, line width=1.1pt] (24,0) -- (34,20) -- (39,30);
  \draw[bkGreen!70!black, line width=1.1pt] (12,30) -- (21,20) -- (23,20) -- (37,0);
  \node[font=\sffamily\scriptsize, text=bkBlue!80!black, anchor=west] at (46,30) {ふつう};
  \node[font=\sffamily\scriptsize, text=bkRed!80!black, anchor=south] at (39,31) {\R{特急}{とっ|きゅう}};
\end{tikzpicture}
\end{minipage}
\end{column}

\begin{butsuri}
\par\addvspace{3pt}\textbf{データロガー：\R{自動}{じ|どう}でグラフを\R{作}{つく}る\R{機械}{き|かい}}\par\noindent
\noindent\begin{minipage}[c]{0.6\linewidth}
\R{工場}{こう|じょう}や\R{冷蔵庫}{れい|ぞう|こ}、ビニールハウスでは、\R{温度}{おん|ど}センサー（3-3）をつないだ\textbf{データロガー}という\R{小}{ちい}さな\R{機械}{き|かい}が、1\R{分}{ぷん}ごとなどに\R{温度}{おん|ど}を\R{記録}{き|ろく}します。\bkidx{でーたろがー}{データロガー}\R{記録}{き|ろく}した\R{表}{ひょう}をコンピュータに\R{送}{おく}ると、$x$-$t$ グラフと\R{同}{おな}じように、\R{時間}{じ|かん}を\R{横軸}{よこ|じく}にしたグラフが\R{自動}{じ|どう}でかけます。\R{冷蔵庫}{れい|ぞう|こ}のとびらを\R{開}{ひら}けると、\R{中}{なか}の\R{温度}{おん|ど}が\R{上}{あ}がり、とじるとまた\R{下}{さ}がるようすが、グラフで\R{見}{み}えます。
\end{minipage}\hfill
\begin{minipage}[c]{0.37\linewidth}\centering
\begin{tikzpicture}
\begin{axis}[bkplot, axis lines=left, width=46mm, height=34mm, xmin=0, xmax=30, ymin=0, ymax=12, xtick={0,10,20,30}, ytick={0,4,8}, xlabel={\R{分}{ふん}}, ylabel={℃}]
  \addplot[bkBlue, line width=1.1pt, smooth] coordinates {(0,4) (5,4.1) (9,4) (10,8.5) (11,9) (13,7) (16,5.3) (20,4.4) (25,4.1) (30,4)};
\end{axis}
\end{tikzpicture}
\end{minipage}

\par\addvspace{3pt}\textbf{\R{蒸気表}{じょう|き|ひょう}：\R{表}{ひょう}のあいだを\R{内挿}{ない|そう}で\R{読}{よ}む}\par\noindent
\noindent\begin{minipage}[c]{0.5\linewidth}
\R{水}{みず}は、まわりの\R{圧力}{あつ|りょく}（おす\R{力}{ちから}の\R{強}{つよ}さ）が\R{高}{たか}いほど、\R{高}{たか}い\R{温度}{おん|ど}でふっとうします。ふっとうする\R{温度}{おん|ど}と\R{圧力}{あつ|りょく}の\R{表}{ひょう}を\textbf{\R{蒸気表}{じょう|き|ひょう}}といい、\R{発電所}{はつ|でん|しょ}やボイラーの\R{設計}{せっ|けい}に\R{使}{つか}います。\qty{115}{\degreeCelsius}は\R{表}{ひょう}にありませんが、グラフで110℃と120℃のあいだを\R{読}{よ}むと、\R{約}{やく}\qty{170}{kPa}です（kPa はキロパスカル。\R{圧力}{あつ|りょく}の\R{単位}{たん|い}）。\R{圧力}{あつ|りょく}なべの\R{中}{なか}は\R{約}{やく}\qty{200}{kPa}なので、\R{水}{みず}は\R{約}{やく}\qty{120}{\degreeCelsius}でふっとうし、\R{料理}{りょう|り}が\R{早}{はや}くできます。
\end{minipage}\hfill
\begin{minipage}[c]{0.47\linewidth}\centering
\begin{bktable}{colspec={Q[c,m,wd=12mm] *{6}{Q[c,m,wd=6.2mm]}}, rowsep=0.8pt, colsep=1.2pt, cells={font=\scriptsize}}
\R{温度}{おん|ど}（℃） & 100 & 110 & 120 & 130 & 140 & 150 \\
\R{圧力}{あつ|りょく}（kPa） & 101 & 143 & 199 & 270 & 362 & 476 \\
\end{bktable}
\par\vspace{2pt}
\begin{tikzpicture}
\begin{axis}[bkplot, axis lines=left, width=54mm, height=38mm, xmin=100, xmax=152, ymin=50, ymax=500, xtick={100,110,...,150}, ytick={100,200,300,400}, xlabel={℃}, ylabel={kPa}]
  \addplot[bkBlue, line width=1.1pt, smooth] coordinates {(100,101.4) (110,143.4) (120,198.7) (130,270.3) (140,361.5) (150,476.2)};
  \addplot[only marks, mark=*, mark size=1.2pt, bkInk] coordinates {(100,101.4) (110,143.4) (120,198.7) (130,270.3) (140,361.5) (150,476.2)};
  \draw[bkRed, line width=0.7pt, dash pattern=on 2pt off 1.5pt] (axis cs:115,50) -- (axis cs:115,169) -- (axis cs:100,169);
\end{axis}
\end{tikzpicture}
\end{minipage}

\end{butsuri}


%% ---------------- 章末問題 ----------------
\bkendproblems
\bkproblemgroup{\R{基本}{き|ほん}}{1}

\bkproblem \R{水}{みず}を\R{温}{あたた}めて、1\R{分}{ぷん}ごとに\R{温度}{おん|ど}をはかりました。\R{方眼}{ほう|がん}は\R{横}{よこ}10マス、たて10マスです。\R{横軸}{よこ|じく}とたて\R{軸}{じく}に\R{何}{なに}をとり、1マスをいくつにすればよいですか。グラフもかきましょう（ノートに）。\par
{\small \R{時間}{じ|かん}（\R{分}{ふん}）0、1、2、3、4、5\quad \R{温度}{おん|ど}（℃）15.0、18.1、20.9、24.0、27.1、29.9}
\bkpans{\R{横軸}{よこ|じく}は\R{時間}{じ|かん} $t$（\R{分}{ふん}）で1マス0.5\R{分}{ふん}。たて\R{軸}{じく}は\R{温度}{おん|ど}（℃）。0からだと1マス5℃、10℃より\R{下}{した}を\R{省}{はぶ}くと1マス2℃で10〜30℃（こちらが\R{大}{おお}きくかける）。\R{点}{てん}がまっすぐに\R{並}{なら}ぶので、\R{直線}{ちょく|せん}を\R{引}{ひ}く}

\bkproblem グラフは、ミナさんが\R{家}{いえ}から\R{本屋}{ほん|や}へ\R{行}{い}き、\R{家}{いえ}へ\R{帰}{かえ}ったときの $x$-$t$ グラフです。（1）\R{本屋}{ほん|や}は\R{家}{いえ}から\R{何}{なん}mですか。（2）\R{本屋}{ほん|や}に\R{何分}{なん|ぷん}いましたか。（3）\R{行}{い}きと\R{帰}{かえ}りの\R{速}{はや}さは、それぞれ1\R{分}{ぷん}に\R{何}{なん}mですか。
\begin{bkfigure}
\begin{tikzpicture}
\begin{axis}[bkplot, width=72mm, height=40mm, xmin=0, xmax=18.5, ymin=0, ymax=500, xtick={0,1,...,17}, xticklabels={0,,,,4,,,,8,,,,12,,,,16,}, ytick={0,100,...,400}, xlabel={$t$（\R{分}{ふん}）}, ylabel={$x$（m）}]
  \addplot[bkBlue, line width=1.3pt] coordinates {(0,0) (4,400) (9,400) (17,0)};
\end{axis}
\end{tikzpicture}
\end{bkfigure}
\bkpans{（1）\nobreak \qty{400}{m}\quad（2）\nobreak 5\R{分}{ふん}\quad（3）\nobreak \R{行}{い}き $400\div4=100$ で\qty{100}{m}、\R{帰}{かえ}り $400\div8=50$ で\qty{50}{m}}

\bkproblem ある\R{台車}{だい|しゃ}の\R{速}{はや}さは、0〜3\R{秒}{びょう}が\qty{4}{m/s}、3〜5\R{秒}{びょう}は\R{止}{と}まって0、5〜9\R{秒}{びょう}が\qty{2}{m/s}でした（ノートに $v$-$t$ グラフをかこう）。\R{動}{うご}きはじめの\R{位置}{い|ち}を0として、3\R{秒}{びょう}・5\R{秒}{びょう}・9\R{秒}{びょう}の\R{位置}{い|ち}を\R{求}{もと}めましょう。
\bkpans{3\R{秒}{びょう}は $4\times3=12$ で\qty{12}{m}、5\R{秒}{びょう}も\qty{12}{m}、9\R{秒}{びょう}は $12+2\times4=20$ で\qty{20}{m}}

\bkproblem ばねにおもりをつるして、のびをはかりました。\R{点}{てん}はまっすぐに\R{並}{なら}びました。\\
{\small おもり（g）0、20、40、60\quad のび（cm）0、1.5、3.0、4.5}\\
（1）30\,gのときののびを\R{求}{もと}めましょう。（2）のびが3.75\,cmになるのは\R{何}{なん}gのときですか。（3）この\R{直線}{ちょく|せん}をのばして、1000\,gのときののびを\R{予想}{よ|そう}すると\R{何}{なん}cmになりますか。それは\R{正}{ただ}しいと\R{言}{い}えますか。
\bkpans{（1）\nobreak 20\,gと40\,gのまん\R{中}{なか}なので、$(1.5+3.0)\div2=2.25$ で\qty{2.25}{cm}\quad（2）\nobreak 3.0と4.5のまん\R{中}{なか}なので、40\,gと60\,gのまん\R{中}{なか}の\qty{50}{g}\quad（3）\nobreak 20\,gで1.5\,cmなので、$1.5\times50=75$ で\qty{75}{cm}。\R{正}{ただ}しいとは\R{言}{い}えない（\R{強}{つよ}く\R{引}{ひ}きすぎると\R{比例}{ひ|れい}しなくなる。2-15）}

\bkproblem PさんとQさんの $x$-$t$ グラフです。Pさんは、たて1\R{目}{め}もりが5\,mのグラフで、5\R{秒}{びょう}で4\R{目}{め}もり\R{上}{あ}がりました。Qさんは、たて1\R{目}{め}もりが2\,mのグラフで、5\R{秒}{びょう}で8\R{目}{め}もり\R{上}{あ}がりました。どちらが\R{速}{はや}いですか。
\bkpans{Pさんは $5\times4=20$（m）で $20\div5=4$、Qさんは $2\times8=16$（m）で $16\div5=3.2$。Pさんが\R{速}{はや}い}

\bkproblemgroup{チャレンジ}{2}

\bkproblem \R{姉}{あね}が\R{家}{いえ}から1\R{分}{ぷん}に\qty{50}{m}で\R{歩}{ある}いて\R{出発}{しゅっ|ぱつ}しました。10\R{分}{ぷん}\R{後}{ご}に、\R{妹}{いもうと}が1\R{分}{ぷん}に\qty{150}{m}の\R{自転車}{じ|てん|しゃ}で\R{追}{お}いかけました。\R{駅}{えき}は\R{家}{いえ}から\qty{1200}{m}です。（1）\R{妹}{いもうと}が\R{出発}{しゅっ|ぱつ}したとき、\R{姉}{あね}は\R{家}{いえ}から\R{何}{なん}mですか。（2）ノートに\R{2人}{ふたり}の $x$-$t$ グラフをかき、\R{妹}{いもうと}が\R{追}{お}いつく\R{時刻}{じ|こく}と\R{場所}{ば|しょ}を\R{読}{よ}みとって、\R{計算}{けい|さん}で\R{確}{たし}かめましょう。（3）どちらが\R{何分}{なん|ぷん}\R{早}{はや}く\R{駅}{えき}に\R{着}{つ}きますか。
\bkpans{（1）\nobreak \qty{500}{m}\quad（2）\nobreak 15\R{分}{ふん}、\R{家}{いえ}から\qty{750}{m}（\R{姉}{あね} $50\times15=750$、\R{妹}{いもうと} $150\times5=750$）\quad（3）\nobreak \R{姉}{あね}は $1200\div50=24$ \R{分}{ぷん}、\R{妹}{いもうと}は $10+1200\div150=10+8=18$ \R{分}{ぷん}。\R{妹}{いもうと}が6\R{分}{ぷん}\R{早}{はや}い}

\bkproblem \R{物理}{ぶつ|り}・\R{工学}{こう|がく}ではこう\R{使}{つか}うの\R{蒸気表}{じょう|き|ひょう}を\R{使}{つか}います。（1）\qty{125}{\degreeCelsius}でふっとうする\R{圧力}{あつ|りょく}は\R{約}{やく}\R{何}{なん}kPaですか。\R{表}{ひょう}の\R{値}{あたい}のまん\R{中}{なか}で\R{考}{かんが}えましょう。（2）この\R{表}{ひょう}は100〜150℃しかありません。\R{線}{せん}をのばして\qty{300}{\degreeCelsius}の\R{値}{あたい}を\R{予想}{よ|そう}するのは、なぜ\R{危}{あぶ}ないのですか。
\bkpans{（1）\nobreak 120℃の199と130℃の270のまん\R{中}{なか}で、$(199+270)\div2=234.5$、\R{約}{やく}\qty{230}{kPa}（くわしい\R{値}{あたい}は\R{約}{やく}\qty{232}{kPa}）\quad（2）\nobreak はかった\R{範囲}{はん|い}の\R{外}{そと}では、きまりが\R{変}{か}わるかもしれないから（\R{線}{せん}は\R{上}{うえ}へだんだん\R{急}{きゅう}に\R{曲}{ま}がっている）}

\begin{furikaeri}
\begin{checklist}
\item \R{軸}{じく}と\R{目}{め}もりを\R{決}{き}めて、はかった\R{点}{てん}からグラフをかける。\bkfurref{bk:3-4:jiku}
\item $x$-$t$ グラフと $v$-$t$ グラフを\R{読}{よ}み\R{分}{わ}けられる。\bkfurref{bk:3-4:xt}
\item \R{表}{ひょう}にない\R{値}{あたい}を\R{線}{せん}から\R{読}{よ}み、\R{外挿}{がい|そう}の\R{注意}{ちゅう|い}を\R{説明}{せつ|めい}できる。\bkfurref{bk:3-4:naisou}
\end{checklist}
\end{furikaeri}

\bkfukidashi{sora}{$x$-$t$ グラフは、\R{急}{きゅう}なほど\R{速}{はや}かったね。\R{次}{つぎ}の\R{章}{しょう}では、その「\R{急}{きゅう}さ」を\R{数}{かず}で\R{言}{い}う\R{方法}{ほう|ほう}を\R{学}{まな}ぶよ。}
